% PM 트랙 Day 4 슬라이드 본문 — 손으로 쓴 정본(한 프레임 = 한 쪽). Day3_슬라이드_body.tex 와 같은 형식.
% 래퍼: Day4_슬라이드.tex(슬라이드판) / Day4_슬라이드_노트.tex(노트판, 우측에 \note).
% 화면에는 결론만, 설명·근거·예외는 각 프레임의 \note{}에.
% 그림은 ../그림/PM_Day4/*.pdf (벡터). 그림 번호 = 파일명 번호(그림 2-1 = fig_2_1_*). 모든 숫자는 Day4 워크북 완성 예시본(정본)에서.
% 데이터 일자 (a) 2027-02-05 런북·SAC 확정 → 실행 기록 02-28 / (b) 2027-05-28 G4. 컷오버 창 02-26 20:00~02-28 06:00(34h) · PoNR 01:40 · 롤백 최종 시한 16:00 · 트리거 4 발동 0 · 최종 AC 151.95 / VAC −5.13 / CPI 0.966.


\newcolumntype{L}[1]{>{\raggedright\arraybackslash}p{#1}}
\newcommand{\ra}{$\rightarrow$}
\newcommand{\til}{\textasciitilde}
\newcommand{\keymsg}[1]{\par\vfill{\color{mDarkTeal}\bfseries #1}\par}
\newcommand{\figcap}[1]{\par\smallskip{\footnotesize\color{black!60}#1}\par}
\newcommand{\fig}[2][0.66]{\centering\includegraphics[height=#1\textheight,width=\textwidth,keepaspectratio]{#2}\par}
\newcommand{\subhead}[1]{{\color{mDarkTeal}\bfseries #1}\par}
\newenvironment{tbl}[2][\small]{\begin{center}\usebeamercolor[fg]{normal text}\setlength{\tabcolsep}{4pt}\renewcommand{\arraystretch}{1.15}#1\begin{tabular}{#2}\toprule}{\bottomrule\end{tabular}\end{center}}

% =====================================================================
% 도입
% =====================================================================
\begin{frame}{PM 트랙 Day 4 — 오픈은 이벤트가 아니라 이관이다}
\begin{itemize}
\item 4일 특강 · 4일차 · 전환과 종료(Transition · Closing)
\item 사례는 같다 — 온담F\&B홀딩스 ONE ONDAM \textbf{P1}, BL-02 PMB 146.82 / 집행 한도 156.0
\item 시점이 둘 — 런북 확정 \textbf{2027-02-05} / 종료 \textbf{2027-05-28 G4}
\item 1차 컷오버 02-26 20:00\til{}02-28 06:00(34h), PoNR 01:40, 트리거 4 발동 0
\item 오늘의 문서는 프로젝트가 아닌 것을 향해 쓴다 — 결정 요청 칸이 아니라 \textbf{인수자 서명 칸}
\end{itemize}
\keymsg{오늘의 질문 — 오픈 다음 날 아침, 운영이 "우리는 인수한 적 없다"고 말하면 무엇이 없었던 것인가}
\note{표지. 교재 서문. Day3 워크북 완성 예시본(EX-01 대안 1 채택 10-16, MG3 02-19 회복)을 읽어 왔는지 확인한다. Day1은 결정하는 문서, Day2는 분해하고 숫자로 만드는 문서, Day3는 편차와 결정 요청의 문서였다면, Day4는 프로젝트 바깥의 사람이 읽고 서명하는 문서다 — 런북은 운영이 오픈 다음 날 읽고, 종료 보고서는 내부감사가 읽고, 편익 원장은 P1이 사라진 뒤 재무본부가 쓴다. 시점이 둘이라는 점(계획 / 실적 열 분리)을 워크북 §0.1 규칙 3으로 미리 못 박는다. 예상 소요 3분.}
\end{frame}

\begin{frame}{오늘의 흐름}
\begin{columns}[T]
\begin{column}{0.47\textwidth}
\subhead{오전 — 이론 (제1\til{}3장)}
\begin{itemize}
\item \textbf{1교시} 90분 · 전환·종료론과 컷오버 — 세 개의 시계 · 런북 · 게이트 · PoNR · 롤백 · 리허설 — 제1\til{}2장
\item \textbf{2교시} 90분 · 서비스 전환 — ITIL 판본 · SAC · ELS · SLA 3층 · error budget · 72시간 — 제3장
\item \textbf{3교시} 90분 · 종료 — 원가 3층 · 계약 종결 · 지체상금 0 · 이관 10건 — 제4장
\end{itemize}
\end{column}
\begin{column}{0.51\textwidth}
\subhead{오후 — 실습과 교훈·편익}
\begin{itemize}
\item \textbf{4교시} 90분 · \textbf{실습 ①} 컷오버 런북·SAC + 이관 체크리스트·SLA/OLA/UC — 워크북 §1\til{}2
\item \textbf{5교시} 75분 · 교훈과 편익 — Syllk · 여섯 칸 · ITO · 분모 정치 · 손익분기 66\% — 제5장
\item \textbf{6교시} 75분 · \textbf{실습 ②} 종료 보고서·교훈 등록부·편익 원장 — 워크북 §3\til{}5
\item \textbf{마무리} 30분 · 학술 배경 · 20종의 사슬을 거꾸로 · Product 트랙 예고 — 제6장
\end{itemize}
\end{column}
\end{columns}
\keymsg{오전은 "받는 쪽"의 문서를 읽고, 오후는 "받는 쪽"이 서명할 문서를 쓴다}
\note{교재 목차. 1교시가 두 장(제1·2장)을 묶는 이유는 컷오버가 종료론의 첫 실행이기 때문이다 — 런북의 마지막 행이 D+1 인계 서명이다. 2교시는 ITIL 어휘를 판본 중립으로 고정한 뒤 SAC·ELS·SLA를 다룬다. 3교시 종료 보고서는 "절감"과 "초과"가 동시에 참인 원가 3층이 핵심이다. 실습 ①은 런북 항목 3\til{}6(게이트·PoNR·트리거·검증)을 손으로 채우는 것이 중심이고, 실습 ②는 종료 보고서 항목 4·6과 편익 원장 측정 책임자 칸이다. 예상 소요 3분.}
\end{frame}

\begin{frame}{학습 목표}
\begin{columns}[T]
\begin{column}{0.49\textwidth}
\subhead{오전 — 읽을 수 있다}
\begin{itemize}
\item 프로젝트·운영·편익의 세 시계를 분리해 오픈의 위치를 말한다
\item 게이트를 "되돌리는 비용이 달라지는 곳"에 세우고 판정자를 배치한다
\item PoNR과 기술적 롤백 시한을 구별하고 결정권 상향 구간을 정한다
\item SAC·ELS 종료 조건·SLA 3층의 결손을 찾는다
\item 종료 보고서의 숫자를 성과보고·변경 로그의 행으로 추적한다
\end{itemize}
\end{column}
\begin{column}{0.49\textwidth}
\subhead{오후 — 쓸 수 있다}
\begin{itemize}
\item ⑯ 런북 게이트 3·PoNR·트리거 4·검증 6종의 조건을 관측값으로 적는다
\item ⑰ 이관 체크리스트의 인수 서명 4열과 SLA/OLA/UC 3층을 채운다
\item ⑱ 원가 3층(168/156/146.82) 대비 최종 AC 151.95를 분리해 적는다
\item ⑲ 교훈 6칸 중 "적용 대상·소유자" 칸을 실명으로 채운다
\item ⑳ 편익 원장의 측정 책임자·분모·손익분기를 서명란과 함께 만든다
\end{itemize}
\end{column}
\end{columns}
\keymsg{목표의 공통 척도 — 프로젝트가 사라진 뒤에도 그 문서가 누군가의 손에 남아 있는가}
\note{교재 서문의 "세 가지를 미리 말해 둔다"를 목표로 바꾼 것. 첫째, 오늘의 시나리오는 Day3 결정 요청에 스폰서가 답한 뒤의 세계다(EX-01 대안 1, R-19 0.60 배정, 대안 3 미발동). 둘째, 기준선 값을 바꾸지 말고 인용하라 — 종료에서도. 승인 예산 168 대비 16.05를 절감으로 적고 싶은 유혹이 가장 큰 시점이 종료다. 셋째, 모든 문서는 인수자 서명으로 끝난다 — 박준영 07:20, 네 사람의 인수 서명, 정미란 G4, 프로그램 PMO장 개정 요청 접수, 재경팀장 정본 인계. 예상 소요 3분.}
\end{frame}

\begin{frame}{Day3 통제 산출물 5종 \ra{} Day4 전환·종료 산출물 5종}
\begin{tbl}[\scriptsize]{L{3.0cm}L{5.4cm}L{2.3cm}}
Day4 산출물 & Day3에서 받는 입력 & 교재 \\
\midrule
⑯ 런북 · SAC & ⑫ EX-01 대안 1 · ⑭ 심각도 1·2 미해결 0 · ⑬ R-14(두 번 컷오버) · ⑪ CR-08 조건 ③·CR-11 & 제2장 · 3.2 \\
⑰ 이관 · SLA/OLA/UC · ELS & ⑬ 이슈 미결 6건 · ⑭ 인터페이스 47본 문서 · ⑮ 하도급 조건 ③·계약 변경 4·6·7번 · R-16 & 제3장 \\
⑱ 종료 보고서 · 이관 목록 · 계약 종결 & ⑫ BL-02(146.82)·EAC4 150.86 · ⑪ CR-01\til{}11 · ⑬ 잔여 EMV 7.26 · ⑮ 귀책 분해표 & 제4장 \\
⑲ 교훈 등록부 & ⑫ 조기 신호 5·착시 5 · ⑭ 감리 귀속 이의 · ⑬ IS-02 · AF-1-08 & 제5장 전반 \\
⑳ 편익 실현 원장 & ⑫ 편익 축 G(118억) · ⑬ R-12(두 분모) · ⑪ CR-04·CR-09 · Day1 ① BC §4·§6 & 제5장 후반 \\
\end{tbl}
\keymsg{20종의 사슬 — 런북이 시작이고 편익 원장이 끝이다. 되돌아가는 파선은 종결 기록이지 재기준선이 아니다}
\note{교재 서문의 표와 그림 0-1(마무리에서 그림 6-1로 다시 본다). 다섯 산출물은 서로를 인계한다 — 런북 D+1 판정 회의가 ⑰ 이관의 시작이고, SAC 조건부 1항목이 ⑱ 이관 목록 1번이 되며, 이관 6번(두 분모 판정, 나영선·2027-09 이사회)이 ⑳ 원장 BR-05로 간다. 강조할 것 — 종료 보고서의 모든 숫자는 Day3 문서의 어느 행에서 왔는지가 추적되어야 하고, 추적되지 않는 숫자는 종료 시점에 만들어진 숫자다(4.1절). 예상 소요 4분.}
\end{frame}

% =====================================================================
\section{1교시 (90분) — 전환·종료론과 컷오버}
% =====================================================================
\begin{frame}{이 교시에서 배우는 것}
\begin{itemize}
\item 프로젝트의 끝과 서비스의 시작은 같은 날이 아니다 — 세 개의 시계 (1.1)
\item 표준이 종료를 어떻게 정의하는가 — PRINCE2 5활동 · PMBOK 8 · ISO 21502 (1.2)
\item 임시조직의 끝은 제도로 정한다 — 해산은 사람별 일자로 (1.3)
\item 분야별 종료 게이트의 뼈대는 같다 — 증명서·하자 목록·보증·서명 (1.4)
\item 컷오버 — 런북·게이트 3·PoNR·롤백 트리거 4·검증 6종·리허설 3회 (제2장)
\end{itemize}
\keymsg{한 문장 — 오픈은 프로젝트가 시스템을 놓은 사건이 아니라, 운영이 받았다고 서명한 사건이어야 한다}
\note{교재 제1장·제2장 2.1\til{}2.7, 워크북 §1. 전반 40분은 종료론(제1장), 후반 50분은 컷오버(제2장)이며, 2.8 D$-$2 라이선스 사건은 2교시 첫머리에서 그림 2-6으로 다룬다. 이 교시의 실습 연결은 워크북 §1 항목 2\til{}6(시각표·게이트·PoNR·트리거·검증)이다. 수강생에게 미리 묻는다 — 최근 오픈에서 "운영이 받았다"는 서명은 누가 언제 했는가(제1장 생각해 볼 질문 1). 예상 소요 2분.}
\end{frame}

\begin{frame}{오픈 $\neq$ 종료 — 세 개의 시계}
\fig[0.58]{../그림/PM_Day4/fig_1_1_three_clocks.pdf}
\figcap{그림 1-1. 프로젝트 시계(G3 \ra{} 컷오버 \ra{} 2차 \ra{} ELS 종료 \ra{} G4 05-28)·운영 시계(ELS 착수 서명 07:20 \ra{} CAB \ra{} 하자보수 2028-05-28)·편익 시계(첫 측정 2027-04 \ra{} G5 2028-03-31). 오픈은 겹치는 점 하나다}
\keymsg{오픈은 세 시계가 만나는 점이 아니다 — 프로젝트의 마지막 마일스톤과 운영의 첫 달이 겹치는 곳이고, 편익 시계는 아직 시작하지도 않았다}
\note{교재 1.1. 온담의 프로젝트는 오픈(02-28 06:00) 뒤 석 달을 더 산다 — 2차 컷오버 03-27, ELS 종료 04-23, 운영 이관 착수 04-26, R5 05-07, 감리 G4 05-14, SAC 판정 05-21, 계약 종결 05-26, G4 05-28. 운영은 D+1 07:20 박준영의 ELS 착수 서명부터 시작되었고, 편익은 B-08 첫 측정(2027-04, 13.6시간)만 있고 판정은 G5다. "이관"의 정의 — 주는 쪽과 받는 쪽이 있고 받는 쪽이 받았다고 서명해야 끝난다. 런북 마지막 행이 "06:00 오픈"이 아니라 "D+1 07:20 서명"인 이유. 흔한 오류 — 헌장 종료일 = 컷오버일로 적어 PM의 마지막 석 달이 계약상 존재하지 않게 되는 것. 예상 소요 7분.}
\end{frame}

\begin{frame}[shrink=11]{PRINCE2 7 Closing a Project — 5활동과 End Project Report}
\fig[0.56]{../그림/PM_Day4/fig_1_2_prince2_closing.pdf}
\figcap{그림 1-2. Prepare planned closure · Prepare premature closure · Hand over products · Evaluate the project · Recommend project closure — 4 산출물(End Project Report · Lessons Report · Benefits Management Approach · Follow-on Action Recommendations)이 ⑯\til{}⑳에 대응한다}
\keymsg{두 규율 — 인수는 운영의 공식 수용이지 PM의 선언이 아니다 / premature closure는 실패 처리가 아니라 정규 활동이다}
\note{교재 1.2. PRINCE2가 종료에서 주는 규율 두 가지를 강조한다. 첫째, hand over products의 인수(acceptance)는 운영·유지보수 조직의 공식 수용 — ⑯⑰의 서명이 그것이다. 둘째, premature closure는 정규 활동이며 중단에도 인계·평가·권고가 있다. P1 헌장 §14의 중단 조건(편익 전망 손익분기 66\% 미만, 컷오버 창 2회 연속 이탈)과 손실 확정 절차(재경팀장 30일 내)는 이 활동의 사전 문서화이고, 종료 보고서에는 "미발동"으로 기록이 남는다. 흔한 오류 — End Project Report 목차를 그대로 옮기고 절마다 한 문단씩 채우는 것. 표준은 무엇이 있어야 하는가를 말하지 그 값이 어디서 오는가를 말하지 않는다. 예상 소요 7분.}
\end{frame}

\begin{frame}{PMBOK 8판 Closing Focus Area · ISO 21502 §7}
\begin{itemize}
\item PMBOK 8 Closing — 인도물 이전·계약 종결·재무 정산·팀 해산·기록·성과 평가·교훈·\textbf{편익 계획 이전}
\item 6판 4.7 한 프로세스 \ra{} 8판은 계약 종결·재무 정산을 Closing 안에 명시
\item 가장 무거운 문장 — "편익 실현 계획을 프로젝트 후 조직에 이전한다" \ra{} ⑳ 소유 재무본부
\item 운영 준비도(operational readiness)를 종료 조건으로 적는다 \ra{} SAC·⑰
\item ISO 21502 §7 — 종료 조건 확인·인계·계약 종결·기록·교훈·편익 책임 명시. 계약이 인용할 때 검수 조건의 해석 기준
\end{itemize}
\keymsg{세 표준의 공통분모 — "누구에게 넘기는가"가 적혀야 종료다. 편익 계획도, 미해결 사항도, 잔여 리스크도}
\note{교재 1.2. PMBOK 6판의 편익 관리 계획(1.2.6.2)에는 "누구에게 넘기는가"가 없었고 8판은 넘긴다고 쓴다 — 온담 편익 원장의 "소유 재무본부(재경팀장), P1 종료 후 정본 관리자" 한 줄이 그 실행이다. ISO 21502가 실무에서 중요한 순간은 계약서가 "ISO 21502에 따른 프로젝트 관리"를 인용할 때이며, 그때 §7이 검수 조건의 해석 기준이 된다 — 원문 대조를 전제로 말한다(판본은 검색으로 확인). 2026년은 PMBOK 8·ITIL v5·개인정보보호법 개정의 3중 전환기이며, 판본 이름을 외우지 말고 어휘를 판본 중립으로 쓰라는 것이 2교시 3.1의 태도다. 예상 소요 6분.}
\end{frame}

\begin{frame}{분야별 종료 게이트 — 이름은 다르고 뼈대는 같다}
\begin{tbl}[\scriptsize]{L{1.5cm}L{3.6cm}L{2.2cm}L{1.9cm}L{1.6cm}}
분야 & 단계적 증명서 & PoNR / go 조건 & 하자·보증 & 받는 쪽 \\
\midrule
건설·EPC & MC \ra{} 시운전 \ra{} 성능 시험 \ra{} PAC \ra{} FAC & 성능 시험 통과 & punch list A/B/C · DLP & 발주자 \\
SI & 릴리스 검수 \ra{} 최종 검수 \ra{} 잔금·유보금 & 데이터 정본 전환 & 하자보수 1년 · 유지보수 별도 & 운영 조직 \\
제조 NPD & 파일럿 런 \ra{} Run@Rate \ra{} 양산 이관 & 수율·불량률 & ECN backlog & 생산 \\
제약·임상 & 데이터 실사 \ra{} 쿼리 해소 \ra{} DB lock & \textbf{DB lock(규제가 명시)} & — & 규제 제출 \\
\end{tbl}
\keymsg{어휘가 아니라 구조를 빌린다 — 조건부 인수의 정규화(EPC) · PoNR의 명시적 정의(제약) · go 조건의 숫자화(NPD)}
\note{교재 1.4. 뼈대는 다섯 — 단계적 증명서, 하자 목록, 성능 확인, 보증 기간, 받는 쪽의 서명. SAC에 해당하는 것은 EPC의 PAC 조건 목록이고, 미해결 사항 이관 목록은 punch list B(인수 후 기한 내 해소)다. SI가 EPC와 다른 점은 컷오버가 데이터·인터페이스·단말의 세 갈래로 동시에 일어난다는 것이며, 지체상금은 기준선 완료일 대비 지연 일수이므로 Day3의 귀책 분해표가 매월 누적되어야 종료 보고서 한 줄이 된다(3교시 4.3). 제약의 교훈 — 되돌릴 수 없는 시각을 규제가 정의한다. SI의 PoNR은 대개 암묵적이고, 그것이 2.3절의 주제다. 흔한 오류 — SI 계약에 FAC, EPC 계약에 하자보수를 쓰는 것. 예상 소요 6분.}
\end{frame}

\begin{frame}{임시조직의 끝 — Day1의 Lundin \& Söderholm을 회수한다}
\begin{itemize}
\item 전환(Transition) — 이후 상태는 영구조직의 것 \ra{} 전환 완료의 판정권도 영구조직(SAC = 박준영)
\item 제도화된 종료 — 종료 보고서·해산·복귀가 제도로 정해져 있다
\item 두 실무 현상 — 종료 직전 인력 이탈(IS-04·IS-11) · 교훈의 실패(조직이 해산하면 배운 것도 해산)
\item Staw의 처방 — 끝났다고 선언하는 사람과 받았다고 서명하는 사람은 달라야 한다
\item 해산은 사람별 일자로 — 발주사 14명 중 12명 06-01 복귀·2명 P5 / 수행사 05-31 철수·하자보수 4명 비상주
\end{itemize}
\keymsg{"05-31 해산" 한 줄이면 6월 첫 주의 하자 문의는 아무에게도 가지 않는다}
\note{교재 1.3. Day1 1.3절의 네 구조 개념(시간·과업·팀·전환)·네 행위 개념 중 전환과 제도화된 종료를 회수한다. Staw(1976)는 Day3 6.6에서 계속 투자의 처방으로 빌렸다 — 중단의 종료에서는 착수 결정자(김선호·프로그램 이사회)가 아닌 판정자(재경팀장 손실 확정, 프리모템 F3)가, 정상 종료에서는 착수 결정자가 "끝났다"를 선언하고 싶어 하므로 받는 쪽(박준영)의 서명이 판정이 된다. 온담의 OD-POS 2.1 종료 세션(03-19, 박세연 CIO 주관, 15년 기록 아카이브)은 감상이 아니라 Bridges의 "끝냄" 처리이며 2교시 3.9에서 다시 본다. P1 PMO 3명 중 1명 하자보수기 잔류는 "승인 전"으로 이관 목록에 남는다. 예상 소요 6분.}
\end{frame}

\begin{frame}[plain]{}
\fig[0.86]{../그림/PM_Day4/fig_2_1_runbook_timeline.pdf}
\figcap{그림 2-1. 1차 컷오버 런북 시각표 T$-$24h\til{}D+1 — 네 갈래(C·D·I·S) 38행, 게이트 1·2·3 계획/실제(+5·+38·+30), PoNR 01:40, 롤백 최종 시한 16:00, 오픈 판정 04:30 \ra{} 오픈 06:00 \ra{} D+1 07:20 서명. 여유 2h 20m 중 1h 10m 소진}
\keymsg{런북은 시각(時刻)이 주어인 유일한 프로젝트 문서 — 확인 열은 "정상 확인"이 아니라 관측값(건수·해시·성공률)으로}
\note{교재 2.1, 워크북 §1 항목 2. T0 = 02-26(금) 20:00 거래 정지, 창 34h(P-35, 허용범위 0). T$-$24h가 있는 이유 — 전날의 최종 상태 점검·전 백업·컨트롤룸 개설·외부 통보(여덟 행)가 담당자의 체크리스트에만 있으면 그 담당자의 병가로 창 시작 전에 이미 늦는다. 행 7(임시 라이선스 유효 확인)은 D$-$2 사건이 없었다면 없었을 행이다. 갈래를 나누는 것은 병렬 작업을 시각표 위에서 보이게 하고 부분 롤백(RB-3)을 정의하기 위해서다. 실행 기록 열 — 편차 총계 게이트 2 +38·게이트 3 +30·이력 적재 $-$25·배포 +10·오픈 0, 여유 소진 1h 10m. 이 숫자가 2차(03-27)·R5·P2 런북이 이번 런북을 베낄 수 있게 한다. 전체 폭 한 장이므로 상단·하단을 순서대로 짚는다. 예상 소요 8분.}
\end{frame}

\begin{frame}{go/no-go 3게이트와 PoNR 01:40 — 되돌리는 비용이 달라지는 곳}
\begin{tbl}[\scriptsize]{L{1.9cm}L{1.9cm}L{3.9cm}L{1.6cm}L{1.6cm}}
게이트 & 계획 / 실제 & 조건(전부 충족) & 판정자 & 되돌리는 비용 \\
\midrule
1 착수 & 19:30 / 19:35 & 인력 97명·백업 해시·임시 라이선스 유효·CAB 잠금 & P1 PM(위임) & 0 \\
2 정본 전환 & 01:00 / \textbf{01:38} & 마스터 검증 6종 통과·재실행 2회 이내 & \textbf{박세연} & 시간(재기동 2h) \\
\textbf{PoNR} & 01:40 / 01:40 & 신 원장 '정본' 활성화 & — & 시간 + 손실 \\
3 매장 배포 & 12:00 / \textbf{12:30} & 이력 6종·IF 47/47·스모크·RB-2 미발동 & \textbf{정미란} & + 매장 재방문 2.1억 \\
기술적 시한 & 16:00 & 오픈 06:00 $-$ 14h(복원 8h+검증 2h+캐시 4h) & 정미란 & 불가 \\
\end{tbl}
\keymsg{세 게이트를 모두 스폰서에게 맡기면 새벽 1시의 게이트 2는 "스폰서 깨우기 vs 그냥 진행"이 된다 — 데이터는 CIO, 돈과 매장은 스폰서}
\note{교재 2.2\til{}2.3, 워크북 §1 항목 3·4. 게이트의 세 요소 — 조건·판정자·시간 한도. 게이트 2의 38분이 핵심 — 00:30 검증 1차에서 표본 2,000행 중 1건 불일치(가격 마스터 VAT 소수 자릿수). 통과 기준 0과 재실행 허용 2회가 미리 적혀 있었기 때문에 새벽 1시에 토론하지 않고 재실행했다. 게이트 3의 30분은 스폰서의 열람 시간 — "스폰서가 늦었다"가 아니라 "판정자 열람 30분을 계획에 포함"으로 적어야 다음 런북이 배운다. PoNR은 기술적으로 되돌릴 수 없는 시각이 아니라 되돌리면 잃는 것이 생기는 시각 — 01:40부터 이력 적재·03:00 3PL 폴백 배치 출고 지시·매장 캐시가 신 원장에 기록된다. 01:40\til{}16:00의 14h 20m이 "되돌릴 수는 있으나 잃는 구간"이고 롤백 결정권이 PM에서 정미란(부재 시 김선호)으로 한 층 올라간다. 게이트 2 직후에 PoNR을 둔 이유 — 떼어 놓으면 03:00 폴백 배치가 구 원장 기준으로 나가 이중 원장이 된다. 예상 소요 9분.}
\end{frame}

\begin{frame}{롤백 트리거 4 · 롤백의 비대칭 — 롤백도 프로젝트다}
\fig[0.50]{../그림/PM_Day4/fig_2_2_reversal_cost.pdf}
\figcap{그림 2-2. 되돌리는 비용의 계단 — 게이트 1 전 0 \ra{} 시간만 \ra{} 시간+손실(PoNR) \ra{} +매장 방문 2.1억(게이트 3) \ra{} 불가(16:00). 트리거 RB-1\til{}4의 위치와 관측값(1/2 · 08:35/11:00 · 9/31점 · 100\%/99.0\%) — 발동 0}
\keymsg{임계값이 숫자가 아니면 판정은 토론이 된다 — 발동하지 않은 트리거의 관측값이 다음 컷오버의 임계값 근거다}
\note{교재 2.4, 워크북 §1 항목 5(그림 2-3이 트리거 비율 막대). 트리거의 요소 — 지표·임계값·관측 시각(런북 어느 행)·판정자·발동 시 조치·복구 소요. RB-1 마스터 검증 미통과(재실행 2회 후, 박세연, 전체 롤백 손실 0, 다음 창 대안 2 03-12) / RB-2 이력 적재 11:00 미완(PM \ra{} PoNR 후 정미란, 손실 있는 롤백 vs 오픈 연기 10:00) / RB-3 배포 실패 31점(5\%) 또는 직영 10점(오정근·박세연, 매장 갈래 부분 롤백) / RB-4 스모크 99.0\% 미만·PG 1본 실패(정미란, 오픈 연기 \ra{} 부분 운영). RB-2 임계값 11:00은 리허설 환산 12:30보다 보수적 — "언제 실패인가"가 아니라 "언제까지 결정하면 나머지 계획이 살아 있는가". RB-3 판정자가 오정근인 이유 — 매장 재방문 6영업일·2.1억은 영업본부의 비용이므로 Senior User가 판정한다. 비대칭 — 컷오버 계획은 38행, 롤백 계획은 대개 한 줄. PoNR 이후 롤백은 신 원장 거래의 역대사·출고 지시 취소·펌웨어 회수가 있어 컷오버보다 어렵다. 예상 소요 8분.}
\end{frame}

\begin{frame}{데이터 검증 6종 — 이 데이터의 정답을 누가 판정하는가}
\begin{tbl}[\scriptsize]{L{2.2cm}L{2.4cm}L{2.0cm}L{2.0cm}L{2.3cm}}
검증 & 분모 & 리허설 1 (11-23) & 본 컷오버 & 서명 \\
\midrule
① 건수 & 1,842 테이블 & 0 & 0 / 0 & 용역사 PM·PMO 품질 \\
② 금액 합계 & 5년 원장 & $-$41,200원 \ra{} 규칙 정정 & 0 / 0 & \textbf{재경팀장} \\
③ 키 유일성 & PK 1,842 & \textbf{1,204}(브랜드 간 중복, R-10) & 0 / 0 & \textbf{박세연} \\
④ 참조 정합성 & FK 3,916 & \textbf{412}(폐점·단종) & 0 / 0 & 박세연 \\
⑤ 코드 매핑 & 4,318 코드 & \textbf{96}(97.8\%) & 0 / 0 & 수행 통합팀·PMO \\
⑥ 표본 대조 & 2,000 / 5,000행 & 18 & \textbf{1 \ra{} 재실행 \ra{} 0} / 0 & PMO 품질·용역사 PM \\
\end{tbl}
\keymsg{분모 · 통과 기준 0 · 서명자 — 셋 중 하나가 없으면 전환 검증은 자기 채점이다. 리허설 1회차의 미통과 3종은 정본 판정(G2)의 입력이었다}
\note{교재 2.5, 워크북 §1 항목 6. 전환 대상 1,842테이블 7.4TB \ra{} CR-11(5년 필터) 5.1TB, 마스터 6종 2,184,310행·이력 3.21억 행. 여섯 각도 — 건수·금액 합계·키 유일성·참조 정합성·코드 매핑 커버리지·표본 필드 대조. 통과 기준이 "1\% 이내"이면 그 1\%가 어느 테이블에 몰려 있는지를 묻지 않게 된다. 리허설 1회차의 1,204·412·96은 전환 도구의 문제가 아니라 원본 데이터의 정본이 정해지지 않은 문제였고, G2(11-27) 마스터 6종 정본 판정으로 사라졌다 — Day3의 데이터 정본 판정이 여기서 회수된다. 본 컷오버 표본 1건은 성격이 다르다(리허설 3회차 통과 후 발생) — 리허설 시점의 통과를 컷오버 시점의 통과로 알지 말 것. 서명 열 — 재경팀장이 금액 합계에 서명하지 않으면 첫 월마감의 "합계가 다르다"가 전환 오류인지 운영 오류인지 가릴 수 없다. 예상 소요 8분.}
\end{frame}

\begin{frame}{리허설 3회 · 두 번 컷오버 — 결과가 나빠야 좋은 리허설}
\fig[0.52]{../그림/PM_Day4/fig_2_4_rehearsals.pdf}
\figcap{그림 2-4. 1회차 11-23 41h 20m(7.4TB 전체, 창 초과 7h 20m, 미통과 3종) \ra{} 2회차 01-08 33h 10m(5.1TB, 표본 3건) \ra{} 3회차 01-22 31h 40m(롤백 리허설 7h 50m 병행) \ra{} 본 컷오버 32h 50m 상당 / 창 34h}
\keymsg{회차마다 목적이 다르다 — 1회 무엇이 깨지는가 · 2회 창 안에 들어오는가 · 3회 여유와 PoNR 이후 롤백의 실측. 리허설 없이 정한 PoNR은 추측이다}
\note{교재 2.6\til{}2.7, 워크북 §1 항목 7. 리허설의 목적은 통과가 아니라 실측 — 소요 시간, 실패하는 검증, 병목 갈래, 교대 인수 누락. 1회차를 전체 볼륨으로 돌린 덕에 "필터 없이는 창에 들어오지 않는다"가 실측되어 CR-11의 근거가 되었다. 축소 볼륨 리허설의 함정 — 인덱스·락·I/O는 비선형이라 10\%로 3시간이면 전체 30시간이 아니다. 3회차의 롤백 병행(7h 50m, 대사 1,140건)이 기술적 시한 16:00과 결정권 상향의 근거다. 두 번 컷오버(R-14) — 1차 매장·주문 34h / 2차 물류 03-27 22:00\til{}06:00 8h, 나눈 기준은 조직이 아니라 P2 설비 준비 상태(리드타임 38주). 비용 셋 — 임시 상태 4주(매장 = 신 원장, 센터 = 구 SAP, 폴백 배치 2회)·"하지 말 것" 2줄·2차 게이트 1 조건 4개. 1차 하이퍼케어 심각도 2 이상 ≤ 5(실제 4)가 2차 게이트 1 조건 ①이었던 것이 L-15 — 두 컷오버를 달력이 아니라 안정화 지표로 묶는다(그림 2-5). 예상 소요 7분.}
\end{frame}

\begin{frame}{1교시 핵심 정리}
\begin{itemize}
\item 세 개의 시계 — 오픈은 프로젝트의 마지막 마일스톤과 운영의 첫 달이 겹치는 점, 편익 시계는 아직 전
\item 이관 = 받는 쪽의 서명 — PRINCE2 hand over · PMBOK 8 편익 계획 이전 · 해산은 사람별 일자로
\item 게이트는 되돌리는 비용이 달라지는 곳에 — 조건·판정자·시간 한도, 데이터는 CIO·돈은 스폰서
\item PoNR 01:40 $\neq$ 기술적 시한 16:00 — 그 사이 14h 20m의 롤백 결정권은 한 층 위
\item 트리거·검증은 숫자와 서명자로 — 발동 0의 관측값과 리허설 1회차의 실패가 다음 컷오버의 자산
\end{itemize}
\keymsg{컷오버는 되돌릴 수 없는 시각을 정하는 기술이고, 종료는 받는 쪽이 서명하는 문서를 쓰는 기술이다}
\note{교재 제1장·제2장 핵심 정리. 세 오류를 세 문서의 부재로 다시 말한다 — 오픈 = 종료(SAC 부재), 보고서 없는 해산(증명 부재, 2019 스마트오더 세 줄), 소멸하는 미해결 사항(이관 목록 부재, 소유자 실명·기한·승인 여부). 종료에서 양식의 부재는 프로젝트 안에서 발견되지 않는다 — 프로젝트가 사라진 뒤 운영·내부감사·다음 PM이 발견한다. 2교시로 넘어가며 — D$-$2 라이선스 사건(2.8, 그림 2-6)이 게이트 1 조건에 행 하나를 추가하게 만든 인과 사슬과, ITIL 어휘를 판본 중립으로 고정하는 이유(3.1)부터 시작한다. 휴식 전 워크북 §1 항목 3·4·5의 빈 템플릿을 펴 두라고 안내한다. 예상 소요 3분.}
\end{frame}

\section{2교시 (90분) — 서비스 전환: ITIL · SAC · SLA 세 층 · error budget}
% =====================================================================
% 2교시 — 교재 제3장, 워크북 §2(⑰ 서비스 전환 패키지). 모든 숫자는 Day4 워크북 완성 예시본(정본).
% =====================================================================
\begin{frame}{2교시 — 이 교시에서 배우는 것}
\begin{itemize}
\item 약속(SLA)은 그것을 받치는 약속(OLA·UC)의 합보다 강할 수 없다
\item SAC는 거부권이 실재할 때만 장치가 된다 — 판정자 박준영
\item ELS는 기간·종료 기준·인력 축소가 연동된 8주
\item 99.5\%는 219.2분이고, 그 예산이 릴리스 횟수를 정한다
\item 인시던트는 복구로, 문제는 원인 제거로 닫는다 — KEDB 6건
\end{itemize}
\keymsg{2교시의 한 문장 — 약속과 그 약속을 받치는 것을 분리하지 않는 것이 이 장의 모든 오류다}
\note{교재 제3장 도입과 3.11절. 1교시가 "되돌릴 수 없는 시각"(컷오버)이었다면 2교시는 그 다음 날부터 운영 조직이 단독으로 서비스를 지원할 수 있을 때까지의 구간이다. 워크북 §2 ⑰ 서비스 전환 패키지(SAC·ELS·SLA/OLA/UC 3층·error budget policy·KEDB)가 이 교시의 산출물이다. 순서는 어휘(ITIL 판본 8분) → 인수(SAC) → 기간(ELS) → 회의체(CAB) → 약속의 산술(SLA 3층·error budget) → 사건(KEDB·개인정보) → 사람(채택) → 절차의 부작용(D$-$2)이다. 예상 소요 3분.}
\end{frame}

\begin{frame}{ITIL 판본을 정직하게 — v5가 현행, 이 교재는 practice 어휘로 고정}
\fig[0.58]{../그림/PM_Day4/fig_3_1_itil_vocab.pdf}
\figcap{그림 3-1. ITIL Version 5(2026-02-12) 상위 구조 — ITIL Value System · Product and Service Lifecycle 8활동. 이 장은 상위 구조 명칭을 쓰지 않고 practice 명칭 9개만 쓴다}
\keymsg{"ITIL 4가 최신"은 2026-02-12 이후 틀린 문장 — 바뀐 것은 상위 구조·6\ra{}8활동, practice 이름은 그대로}
\note{교재 3.1절. 판본 논쟁을 8분 안에 끝내기 위한 프레임이다. 사실 — 2026-02-12 PeopleCert가 ITIL v5 Foundation 출시, 상위 모듈 3\til{}4월 순차. 바뀐 것 둘: Service Value System \ra{} ITIL Value System, Service Value Chain 6활동 \ra{} Product and Service Lifecycle 8활동. 유지된 것: 34 practices·7 원칙·4 차원. 이 교재가 ITIL 4 어휘(Change Enablement·Service Validation and Testing·Release·Deployment·Incident·Problem·Service Level Management·ELS·KEDB)로 고정하는 이유 셋 — 회사 운영 규정과 벤더 계약이 그 어휘 / practice 이름은 판본 무관 / 다투는 것은 어휘가 아니라 판단. 회사에 돌아가 규정을 개정할 때 바꿀 것은 상위 구조 명칭과 6\ra{}8 매핑뿐이다. 서비스 전환(Transition)은 8활동 중 Transition과 Operate의 경계이고 ELS는 그 경계 위의 기간이다. 예상 소요 6분.}
\end{frame}

\begin{frame}{SAC — 인수 거부권이 실재해야 문서가 장치가 된다}
\begin{itemize}
\item 판정자 = 인수자(운영 조직, 박준영) — 프로젝트가 판정하면 자기 평가
\item 거부권은 헌장 §14 ①에 걸려 있어야 실재한다
\item 조건부 충족 경로 — 없으면 1항목 때문에 전체 거부 또는 거짓 충족
\item 시점 — G1 초안 \ra{} 월 검토 \ra{} G3 사전 점검 \ra{} G4 판정
\item 항목 12 · 각 항목에 판정자와 증빙 열
\end{itemize}
\keymsg{프리모템 12번("미해소 23건으로 인수 거부")과의 차이는 항목 수가 아니라 G3 사전 점검의 유무다}
\note{교재 3.2절 개념, 워크북 §2 항목 1\til{}2. SAC(Service Acceptance Criteria)는 ITIL Service Validation and Testing이 원전, PRINCE2 7 Quality practice의 acceptance criteria \ra{} acceptance record 사슬. 항목 10\til{}12개 — 운영 문서·모니터링·백업/복구·보안·교육·Known Error·SLA/OLA/UC 서명·인시던트/문제/변경 절차·소스/형상·인터페이스 운영·라이선스/계약 이관·잔여 결함/ELS 종료. 장치가 되는 조건은 하나 — 판정자가 인수자이고 그에게 거부권이 실재. 거부권이 실재하면 두 번째 조건이 따라온다 — 조건부 충족의 경로("미해소 0 또는 기한·책임자 서명"). 조건부가 없으면 인수자는 1항목 때문에 전체를 거부하거나(하자보수기 지연·유보금·팀 해산 지연) 거짓으로 충족을 서명한다. 예상 소요 6분.}
\end{frame}

\begin{frame}{SAC 12항목 — G3 사전 점검(미해소 7)에서 G4 판정(충족 11·조건부 1)까지}
\fig[0.58]{../그림/PM_Day4/fig_3_2_sac_progress.pdf}
\figcap{그림 3-2. 01-28 미해소 7 \ra{} 14주 해소 이벤트(DR 시험 04-10 · 임계값 12\ra{}24 · 데스크 교육 11/11 · CAB 8회 · UC 재협상 · 문서 12종 05-14) \ra{} 05-21 판정. 조건부 1(항목 11 아카이브 접근 절차, 재경팀 06-15) = 이관 목록 1번}
\keymsg{사전 점검이 있었기 때문에 미해소 7은 작업 목록이 되었고, 조건부 1은 이관 목록 1번이 되었다}
\note{교재 3.2절 온담 사례, 워크북 §2 항목 2·§1.5 항목 8. SAC v1.0은 G1(2026-04-24) 초안 합의(RACI 22행), 월 검토 9회, 사전 점검 2027-01-28, 판정 05-21, 판정권자 박준영. 01-28 미해소 7 — 운영 문서 8/12, 모니터링 임계값 12/24, DR 시험 미실시, 데스크 교육 6/11, SLA UC 3PL 없음·수행사 UC 초안, CAB 미실시, 라이선스·계약 이관 미해소. 조건부 1은 항목 11의 아카이브 접근 절차 — 문서와 운영팀 검토는 끝났으나 재경팀의 세무 보존 요건 정합 확인이 06-15에야 나온다. 이것을 "미해소"로 두면 G4가 밀리고 "충족"으로 두면 거짓이다 — 그래서 조건부 경로가 필요하고, 조건부는 기한(06-15)·책임자 서명을 갖고 이관 목록 1번이 된다. 항목 7이 "SLA 서명"이 아니라 "3층 전부 서명"인 이유는 뒤의 SLA 프레임에서. 예상 소요 6분.}
\end{frame}

\begin{frame}{ELS 8주 — 종료 조건이 없는 하이퍼케어는 상주가 된다}
\fig[0.58]{../그림/PM_Day4/fig_3_3_els_curves.pdf}
\figcap{그림 3-3. 2027-03-01\til{}04-23 · 오류율 2.4\ra{}0.4\% · 콜 1,840\ra{}450 · P1 2(1주차)·P2 9 · 상주 20\ra{}15\ra{}10\ra{}6\ra{}4 · 5주차 재상승은 2차 컷오버 유예 조건 안 · 7·8주 2주 연속 충족 \ra{} 04-23 종료}
\keymsg{ELS의 세 요소 — 기간 · "n주 연속"의 종료 기준 · 지표와 연동된 인력 축소. 하나라도 빠지면 상주는 하자보수가 된다}
\note{교재 3.3절, 워크북 §2 항목 5. 종료 기준 다섯 — ① 2주 연속 P1 0, ② P2 주 2건 이하, ③ 매장 주문 오류율 0.5\% 이하(O-13 1.8\% \ra{} 목표), ④ 데스크 콜 480/주 이하, ⑤ 심각도 2 backlog 0. 인시던트 140 = P1 2·P2 9·P3 41·P4 88, 핫픽스 11, 판정 박준영. 상주 축소가 달력이 아니라 지표에 연동되어 있다 — 3주차 15는 2주차 판정(심각도 2 이상 4 $\le$ 5) 후, 6주차 10은 오류율 0.5\% 도달 후, 8주차 6은 7주차 기준 충족 후. 5주차 재상승(P2 2·오류율 0.8\%, P2-07 폴백 배치 84분)은 "2차 후 1주는 기준 적용 유예"가 미리 적혀 있었기 때문에 ELS 연장 토론이 없었다. 가용성 산정 — 3월 다운 57분은 P1-01만 산입, P1-02(포인트 이중차감)는 데이터 손상이지 다운이 아니다. 산입 규칙이 미리 없으면 "왜 안 넣었나"가 된다. 예상 소요 7분.}
\end{frame}

\begin{frame}{Change Enablement · Release · Deployment — 셋을 구분해야 CAB가 정치가 되지 않는다}
\begin{tbl}[\scriptsize]{L{0.2\textwidth}L{0.24\textwidth}L{0.48\textwidth}}
practice & 답하는 질문 & 요소 \\ \midrule
Change Enablement & 이 변경을 해도 되는가(승인) & standard·normal·emergency · change authority · CAB/ECAB · 변경 일정 \\
Release Management & 무엇을 언제 묶어 내보내는가(구성) & 릴리스 단위 · 릴리스 일정 · 릴리스 노트 \\
Deployment Management & 어떻게 환경에 옮기는가(실행) & big-bang·phased·pull·continuous · 배포 절차 · 롤백(런북) \\
\end{tbl}
\begin{itemize}
\item 섞이면 CAB가 "이 기능을 넣을 것인가"와 "이번 주 배포해도 되는가"를 한 회의에서 다룬다
\item 온담 — CCB(정미란) 05-28 해산 · CAB(박준영, 수요일) 03-17 첫 회, 8회 = SAC 항목 8 증빙
\end{itemize}
\keymsg{무엇을 넣을지는 릴리스 계획(현업), 언제 어떻게 내보낼지는 CAB(운영 위험 · error budget), 표준 변경은 CAB에 오지 않는다}
\note{교재 3.4절, 워크북 §2 항목 6. 세 practice는 다른 질문에 답한다. 섞이면 제품 결정(현업 우세)과 운영 위험 결정(운영 우세)이 한 회의에 올라 정치가 되고, 나누면 계산이 된다 — 그 계산의 환율이 error budget이다. 운영 이관은 변경 권한의 이전이기도 하다 — 프로젝트 중 CCB(정미란 위원장), 이관 후 CAB(박준영 위원장). 겹치는 기간(ELS)의 규칙: CAB 첫 회 03-17\til{}CCB 해산 05-28 사이 프로젝트 변경(CR-14 잔여·R5 릴리스)은 CCB, 운영 변경(핫픽스·파라미터)은 CAB. 이관 3번(22:00 운용 전환)은 CCB 01-21에서 "P2 안정화 후 CAB 재상정"으로 넘어간 것 — 프로젝트 변경이 운영 변경으로 권한 이전된 실물. 표준 변경 목록이 없으면 파라미터 하나도 2주 걸린다. 예상 소요 6분.}
\end{frame}

\begin{frame}{SLA/OLA/UC 3층 — 한 층이 비면 SLA는 협상이 아니라 산술로 불가능하다}
\fig[0.56]{../그림/PM_Day4/fig_3_4_sla_chain.pdf}
\figcap{그림 3-4. 항목 ② P1 복구 2h — 사슬 15+30+30+240 = 315분 $>$ 120분(UC 초안 불가) \ra{} UC 재협상(착수 30·복구 90, 연 +0.6억) 165분 \ra{} 병렬 착수 $\le$ 120 성립. 항목 ⑥ 출고 확정 D+1 06:00 — 3PL UC 없음 \ra{} 최선 노력 강등}
\keymsg{SLA 복구 시간 $\ge$ 받치는 OLA·UC 사슬의 합 — 강등은 패배가 아니라 정직한 경계 표시이자 재협상 안건 목록이다}
\note{교재 3.5절, 워크북 §2 항목 7(3층 매핑표). SLA = IT본부 \ra{} 현업, OLA = IT본부 내부 팀 간(서비스운영 \ra{} 인프라·보안·개발운영), UC = 외부 공급자(수행사 유지보수·CSP·앱 벤더·3PL·API GW 벤더). 온담 SLA 7·OLA 3·UC 5, 서명 05-14. ②는 G3 사전 점검(01-28)에서 315분이 발견되어 05-14 전에 재협상이 끝났다 — G4에서 처음 발견되었다면 SLA는 발효되지 못한 채 이관되었을 것. ⑥은 출고 확정이 대성로지스 WMS의 작업이고 계약은 물류본부 소관(P1 예산 밖), 준실시간 연동은 거부(R-03·IS-05) — UC 자체가 없으므로 어떤 OLA로도 성립하지 않는다. 2027-06 3PL 재협상 후 재상정 = 이관 4번(한동석, 06-30). 나머지: ① 가용성 99.5\%(CSP 99.9\%), ③ P2 8h(330 $\le$ 480), ⑦ 앱 주문 연동 4h(210 $\le$ 240, R-18 월 3,600만 원, 이관 5번). 예상 소요 8분.}
\end{frame}

\begin{frame}{Error budget — 99.5\%를 3시간 39분으로, 다시 릴리스 속도로}
\begin{itemize}
\item 평균 달 = 365.25/12 $\times$ 1,440 = 43,830분 · 0.5\% = \textbf{219.2분 = 3h 39m}
\item 계획 정지(CAB 승인 창·2차 컷오버 8h)는 제외 — 산입하면 3월 98.7\%
\item 비계획 유보 60 : 릴리스 몫 40 = 131.5분 : \textbf{87.7분}
\item 릴리스당 기대 다운 = 변경 실패율 $\times$ 복구 시간(DORA)
\item policy — 유보 초과 시 표준 릴리스 동결 · 릴리스 몫 소진 시 CAB 승인만 · 결정자 박준영
\end{itemize}
\keymsg{같은 99.5\%를 운영은 "지킬 숫자"로, 프로젝트는 "넘지 않을 숫자"로 읽는다 — 예산으로 바꾸면 토론이 계산이 된다}
\note{교재 3.6절 개념, 워크북 §2 항목 3(error budget policy). SRE(Beyer 외, 2016, 3·4장). 28일 달 202분·30일 216분·31일 223분. 예산은 장애(비계획)에만 쓰는 것이 아니라 릴리스에도 쓴다 — 릴리스 한 번은 기대값으로 실패율 $\times$ 복구 시간만큼 먹는다. policy가 없으면 예산이 소진되어도 아무 일도 안 일어나고, 그러면 예산이 아니다. 계획 정지 제외 규칙이 미리 없으면 2차 컷오버 8시간이 들어가 헌장 §14 ④가 흔들린다. ELS 실적: 3월 다운 57분(P1-01) = 223분의 26\%, 유보의 43\% / 4월 0분. R5 릴리스(05-07)와 핫픽스 11건이 이 예산 안에서 결정되었다. 다음 프레임에서 환전 계산. 예상 소요 5분.}
\end{frame}

\begin{frame}[shrink=8]{환전 계산 — 릴리스 속도를 여섯 배로 만드는 것은 기능이 아니라 복구 시간이다}
\fig[0.34]{../그림/PM_Day4/fig_3_5_error_budget.pdf}
\figcap{그림 3-5. 219.2분 = 유보 60\% + 릴리스 40\%(87.7분) · 3월 실측 57분}
\begin{tbl}[\scriptsize]{L{0.34\textwidth}L{0.14\textwidth}L{0.2\textwidth}L{0.2\textwidth}}
릴리스당 기대 다운 & 분 & 월 가능(87.7 $\div$) & 주 환산 \\ \midrule
DORA 실측 12\% $\times$ 240분 & 28.8 & 3.0회 & 0.7회 — 주 1회 불가 \\
중간(하자보수기) 10\% $\times$ 90분 & 9.0 & 9.7회 & 2.2회 — 표준 주 1회 + 핫픽스 \\
ELS 목표 10\% $\times$ 45분 & 4.5 & 19.5회 & 4.5회 \\
\end{tbl}
\keymsg{240분 \ra{} 45분은 KEDB(우회 준비된 오류) · 회귀 자동화(R-19 0.25억) · 롤백 스크립트의 문제 — 안정과 속도의 환율}
\note{교재 3.6절 온담 사례, 워크북 §2 항목 3. DORA 실측(2026-09, Day3 2.7절)으로는 월 3회가 한계이고 ELS 목표로는 19회. 차이를 만드는 변수는 실패율(12\ra{}10\%)보다 복구 시간(240\ra{}45분)이다. 복구 45분은 기능 팀이 아니라 KEDB·회귀 자동화·롤백 스크립트의 문제 — 그래서 KEDB 6건이 SAC 항목 6이고 R-19 회귀 도구 0.25억이 릴리스 속도 투자다. 이 환율이 있으면 CAB는 "위험하다/필요하다"의 정치가 아니라 "이번 릴리스가 몫의 몇 분을 먹는가"의 계산이 된다. 실습에서 수강생은 자기 시나리오의 가용성 목표를 분으로 환전하고 릴리스 몫을 나누어 본다 — 99.9\%면 43.8분, 릴리스 40\%면 17.5분, DORA 실측이면 월 0.6회다. 예상 소요 6분.}
\end{frame}

\begin{frame}{Incident와 Problem, Known Error DB — 복구로 닫는 것과 원인 제거로 닫는 것}
\begin{tbl}[\scriptsize]{L{0.1\textwidth}L{0.3\textwidth}L{0.36\textwidth}L{0.14\textwidth}}
KE & 증상 \ra{} 우회 & 근본 원인 · 상태 & 소유자 \\ \midrule
KE-01 & 구형 단말 9점 인쇄 8초 \ra{} 재부팅 & 펌웨어 호환 — 협의 중(이관 2번) & 오정근 \\
KE-02 & 폴백 배치 겹침(R-16) \ra{} 22:30 & 3PL 준실시간 전환 시(2027-06, 이관 4번) & 한동석·박준영 \\
KE-03 & 포인트 이중차감 \ra{} 중복 검출 배치 & 멱등 키 — 03-12 핫픽스 제거, 감시 잔류 & 수행사 \\
KE-04 & IF-SAP-12 타임아웃 \ra{} 수동 재전송 & ECC 야간 배치 경합 — EoM 2027-12까지 유지(이관 7번) & 박세연 \\
KE-05 & 구형 xls 수량 0 \ra{} xlsx 안내 & 사용률 1\% 미만 — 보류(CAB 04-14) & 박준영 \\
\end{tbl}
\keymsg{인시던트 140 · 문제 9 = 근본 제거 3 + Known Error 6 — 상태가 다르므로 소유자가 다르고, 소유자가 다르므로 이관처가 다르다}
\note{교재 3.7절, 워크북 §2 항목 4(KEDB). Incident Management는 복구, Problem Management는 원인 제거. 인시던트는 원인을 몰라도 복구되면 닫히고, 문제는 재발이 없어도 원인이 제거되어야 닫힌다. Known Error = 원인이 알려졌고 우회가 있으나 근본 제거가 아직인(또는 안 하기로 한) 오류. 표에 없는 KE-06(옴니 매장 간 재고 이동 예약 반영 15분 지연, KE-02와 동일 원인 계열, 오정근)까지 6건. 상태 네 종 — 제거 완료(KE-03), 다른 과제 의존(KE-02·06 \ra{} 3PL, KE-04 \ra{} ECC), 의도적 보류(KE-05), 협의 중(KE-01). Day3 등록부의 R-13(구형 단말)·R-16(배치 충돌)이 KE-01·KE-02가 된 것이 "잔여 리스크의 운영 이관"의 실물(4.6절, 3교시). "보류" 결정을 적지 않으면 6개월 뒤 "왜 안 고쳤나"가 된다. 예상 소요 6분.}
\end{frame}

\begin{frame}{개인정보 사건 — 조문은 단정, 적용은 비단정, 시한은 단정}
\fig[0.52]{../그림/PM_Day4/fig_3_6_privacy_timeline.pdf}
\figcap{그림 3-6. P1-02 포인트 이중차감(2,987명, 03-05 14:20 인지) — CPO 의견 20h40m \ra{} 결정 \ra{} 통지·복구 27h40m \ra{} 신고 49h40m \ra{} 시한 72h. 이 절차가 운영 문서에 없어 회의를 소집한 사실이 L-14}
\keymsg{개정법(제21445호, 2026-09-11 시행) — '유출등'에 훼손 포함·72시간은 단정, 이 사건이 훼손인지는 비단정, 그래도 시한 안에 결정한다}
\note{교재 3.8절, 워크북 §0.1 규칙 7·§2 항목 8. 조문 수준에서 단정 가능한 것: 유출등의 정의에 위조·변조·훼손 포함, 인지로부터 72시간 내 정보주체 통지·보호위원회 신고(대상 요건은 시행령·고시), 과징금 상향. 단정할 수 없는 것: ① 시스템 오류로 잔액이 틀리게 기록된 것이 '훼손'인가(시행 초기 해석례 부재), ② 2,987명이 규모 요건에 해당하는가, ③ 복구 후에도 신고 의무가 유지되는가. 최영주 CPO 서면 의견(03-06 11:00): 배제할 수 없으므로 보수적으로 72시간 내 신고 권고, 통지·복구는 규제 해당 여부와 무관하게 즉시. 결정 03-06 14:00(정미란·최영주·박준영·PM) — 통지·복구 03-06 18:00, 신고 03-07 16:00(49h40m), 협의회 서면 03-08. 기록을 "법상 신고 대상이므로 신고"로 정리하면 깔끔하지만 거짓이다 — 셋을 분리해 적어야 다음 사건에서 "지난번 쟁점 ①이 그 후 결정례로 정리되었는가"를 볼 수 있다. CPO의 의견은 검토 의견이지 판정이 아니다. 예상 소요 8분.}
\end{frame}

\begin{frame}{사용자 채택과 변화관리 — 저항은 정보다}
\begin{tbl}[\scriptsize]{L{0.18\textwidth}L{0.3\textwidth}L{0.12\textwidth}L{0.3\textwidth}}
집단 & 지표(2주 \ra{} 8주) & 장벽 & 조치 \\ \midrule
가맹 점주 388점 & 팩스·엑셀 발주 24\% \ra{} 11\% (목표 $\le$15) & K/A & P3 방문 교육 2회차 · R5 간편 재발주 \\
가맹 점주 & 마감 20:00 준수 81 \ra{} 96\% & R & 마감 30분 전 알림 \\
직영 227점 & 반품 절차 준수 78 \ra{} 94\% & \textbf{R} & POS 강제 확인 단계 + 지역매니저 점검 \\
서비스데스크 11명 & 1차 해결률 52 \ra{} 73\% & K & FAQ 20건 · KEDB 우회 공유 \\
물류센터 246명 & 재고 조회 사용률 87\%(4주) & A/D & 한동석 서면 04-12 \\
\end{tbl}
\keymsg{채택은 교육 이수율이 아니라 사용 지표로 잰다 — R 장벽에 교육을 더 하면 ELS가 끝나는 날 내려간다}
\note{교재 3.9절(그림 3-7), 워크북 §2 항목 9. ADKAR(Prosci, Hiatt 2006) — Awareness·Desire·Knowledge·Ability·Reinforcement, 어느 단계에 막혀 있는지의 진단이 먼저다. Kotter의 단기 성과 공지 — 04-05 협의회 공문 "발주 리드타임 38h \ra{} 13.6h". Bridges의 끝냄 — OD-POS 종료 세션 03-19. Ford·Ford·D'Amelio(2008): 저항은 극복할 장애가 아니라 변화의 설계에 관한 정보이고 때로는 추진자 자신이 만든 것이다 — 강윤지 "이번엔 없어지지 않는 건가요"는 2019년의 기억이 만든 D 장벽, 김미르 "처리량 목표가 인력 절반이라는 뜻이라면"은 편익 지표와 인력 평가의 연결에 대한 A 장벽. 가맹의 채택률이 오른 것은 교육이 아니라 부속서 동의(12-22, 91\%)로 D가 먼저 풀렸기 때문이고, 직영의 절차 미준수(R)가 P1-02의 노출면이었다 — 채택과 인시던트가 같은 자리에서 만난다. 예상 소요 6분.}
\end{frame}

\begin{frame}{D$-$2 라이선스 사건 — 거버넌스 변경의 부작용이 14개월 뒤 청구되는 구조}
\fig[0.55]{../그림/PM_Day4/fig_2_6_license_chain.pdf}
\figcap{그림 2-6. 내부통제 지침(2026-02-16) · 50/50 계약 변경(03-13) · 조달 상태 표 "계약 내" \ra{} 02-13 CFO 결재 대기 \ra{} 02-24 14:30 발견 \ra{} ② 임시 라이선스 0.02억 \ra{} L-09 \ra{} 지침 개정 요청 2·3}
\keymsg{세 개의 정상적인 결정이 겹쳐 만든 대기 — 담당자가 아니라 규정의 겹침을 적어야 L-09가 지침 개정이 된다}
\note{교재 2.8절(1교시에서 넘긴 절), 워크북 §1 항목 9(런북 부록). 2027-02-24(수) 14:30 컨트롤룸 사전 점검에서 API GW 운영 환경 2차분 발주서 "미접수" 회신. 발주서는 02-13 박세연 결재 후 CFO 직접 결재 대기 — 재무본부 내부통제 지침(2026-02-16)이 5,000만 원 이상 단일 건을 CFO 직접 결재로 바꿨고, 정미란은 출장 중. 선택지 ① 결재 대기 \ra{} 03-04, 창 이탈 / ② 벤더 임시 라이선스 30일 0.02억(컨틴전시 미배정 풀, CCB 03-18 사후 보고) / ③ 개발 키 운영 사용 — 계약 위반, 기각. 원인은 실수가 아니라 (1) R-04 대응으로 라이선스를 50/50으로 바꾼 계약 변경(IS-01, 2026-03-13)이 2차분 발주를 2027-02로 밀었고 (2) 그 시점이 지침 시행 후였으며 (3) 조달 상태 표가 "계약 내"로만 표기해 발주 상태가 보이지 않았다는 것. 그래서 런북 부록에 적었고 다음 컷오버 게이트 1 조건에 "라이선스·결재 상태"가 들어가며, L-09 \ra{} 내부통제 지침 개정 요청 2번(게이트·창 앞 4주 이내 발주 건 사전 일괄 결재). 이 사건이 2교시에 있는 이유 — 서비스 전환의 사건 대부분은 기술이 아니라 절차의 겹침에서 온다. 예상 소요 6분.}
\end{frame}

\begin{frame}{2교시 핵심 정리 — 약속과 그 약속을 받치는 것}
\begin{itemize}
\item ITIL 현행은 v5 · 이 교재는 practice 어휘로 고정 — "4가 최신"이라 쓰지 마라
\item SAC = 인수자 판정 + 헌장 거부권 + 조건부 경로 · 미해소 7 \ra{} 충족 11·조건부 1 \ra{} 이관 1번
\item ELS 8주 = 기간 · n주 연속 종료 기준 · 연동된 인력 축소(20\ra{}4)
\item SLA $\ge$ OLA·UC 사슬의 합 — 315 $>$ 120 재협상, 3PL UC 없음 강등
\item 99.5\% = 219.2분 · 릴리스 몫 87.7 · 복구 240분이면 월 3회, 45분이면 19.5회
\end{itemize}
\keymsg{인시던트는 복구로, 문제는 원인 제거로, 개인정보는 조문·적용·시한 세 줄로, 채택은 사용 지표로 — 워크북 §2 ⑰로 간다}
\note{교재 제3장 핵심 정리 5항. 4교시 실습 ①에서 워크북 §2 ⑰ 서비스 전환 패키지를 쓴다 — SAC 12항목의 판정자·증빙 열, ELS 종료 기준 다섯과 인력 축소 연동, 3층 매핑표에서 가장 긴 사슬 합산, error budget policy(환전 계산 포함), KEDB 6건의 상태·소유자·이관처. 수강생에게 미리 말할 것 둘 — 자기 시나리오의 SLA 항목을 하나 골라 3층을 적어 보면 대부분 UC가 없거나 UC가 SLA보다 길다는 것을 발견한다 / 강등은 실패가 아니라 재협상 안건이다. 3교시는 종료(제4장) — 종료 보고서·BL-02 대비 최종 성과·계약 종결·이관 10건. 예상 소요 4분.}
\end{frame}

\section{3교시 (90분) — 종료: 보고서는 마지막 독자를 위해 쓴다}

% =====================================================================
% 3교시 — 종료 (교재 제4장, 워크북 §3 ⑱)
% =====================================================================
\begin{frame}{3교시에서 배우는 것 — 선언이 아니라 증명}
\begin{itemize}
\item 종료 보고서 8항목 — 모든 숫자가 어느 문서 어느 행에서 왔는지 추적된다
\item BL-02 대비 최종 4축 + 리스크·편익·지속가능성, 원가는 3층을 분리해 적는다
\item 계약 종결 — 검수 \ra{} 정산 \ra{} 유보금 \ra{} 하자보수 \ra{} 지체상금 0의 근거
\item 문서·형상 인수(2023 지적 회수) · 이관 10건 · 잔여 리스크 · 팀 해산
\item 마지막 독자 — 3년 뒤의 내부감사·분쟁·다음 PM
\end{itemize}
\keymsg{종료 보고서는 프로젝트의 마지막 문서가 아니라 내부감사의 첫 문서다}
\note{교재 제4장 도입. 2교시(서비스 전환)에서 SAC 판정 05-21·ELS 종료 04-23이 나왔고, 이 교시는 그 둘을 입력으로 받아 G4(2027-05-28) 종료 보고서를 쓴다. 2019년 세 줄짜리 종료 보고서와 2027년 여덟 항목의 차이는 길이가 아니라 "숫자의 출처가 있는가"임을 먼저 못 박는다. 워크북 §3(⑱)의 항목 1\til{}8이 이 교시의 순서와 같다. 예상 소요 4분.}
\end{frame}

\begin{frame}{종료 보고서의 구조 — 8항목과 입력 목록}
\begin{tbl}[\scriptsize]{L{0.06\textwidth}L{0.40\textwidth}L{0.44\textwidth}}
\textbf{\#} & \textbf{항목} & \textbf{입력(어느 문서에서 오는가)} \\ \midrule
1\til{}2 & 문서 정보·완료 기준 판정 / BL-02 대비 최종 성과 & 헌장 §14 다섯 기준 · 성과보고 제17호(05-31) \\
3\til{}4 & 변경 로그 총괄 / 컨틴전시·관리예비비 집행 총괄 & 변경 로그 v1.4 · 결재·CCB 월 보고 \\
5\til{}6 & 잔여 리스크 이관 / 미해결 사항 이관 목록 & 등록부 v1.4 · SAC 판정 · KEDB · SLA \\
7 & 계약 종결 & 계약 종결 서면 7건 · 검수 확인서 R1\til{}R5 \\
8 & 문서·형상 인수 · 팀 해산 · 승인 & 클린 빌드 05-12 · SBOM · 편익 원장 v1.0 \\
\end{tbl}
\begin{itemize}
\item 첫 문장 — 기간 2026-01-05\til{}2027-05-28(17개월), BL-01 \ra{} BL-02, \textbf{재기준선 0회}
\item 완료 기준 다섯 — 각각 판정·근거·일자·\textbf{판정자 실명}(박준영 SAC · 박세연 형상)
\end{itemize}
\keymsg{추적되지 않는 숫자는 종료 시점에 만들어진 숫자다 — 첫 줄은 입력 목록이다}
\note{교재 4.1. PRINCE2 7 End Project Report(PID 대비 성과·편익 리뷰 계획·교훈·후속 조치)와 PMBOK 8판 Closing Focus Area(성과 평가·계약 종결·재무 정산·기록 보관·팀 해산)를 합친 것이 온담의 8항목이다. 증명해야 하는 것 셋은 2020·2023년 지적사항에서 온다 — 편익 산정 근거, 예비비 집행 승인 근거, 소스·형상 인수. "실무에서 틀리는 것" — 종료 시점에 처음 쓴다(성과보고 17호가 있으면 항목 2는 그 복사다), 완료 기준을 종료 시점에 정한다, 입력 목록이 없다. 워크북 §3.3 항목 1. 예상 소요 8분.}
\end{frame}

\begin{frame}{BL-02 대비 최종 성과 — 네 축, 편차 0·0·$-$5.13·24\%}
\begin{tbl}[\scriptsize]{L{0.08\textwidth}L{0.28\textwidth}L{0.30\textwidth}L{0.12\textwidth}L{0.10\textwidth}}
\textbf{축} & \textbf{BL-02} & \textbf{최종 실적(05-28)} & \textbf{편차} & \textbf{판정} \\ \midrule
범위 & FP 12,585 / RTM 1,189 & 인수 FP 12,585(100\%) & 0 & 이내 \\
일정 & G4 2027-05-28 & G4 05-28 정시, 창 34h·8h 정시 & 0영업일 & 이내 \\
원가 & PMB 146.82 & AC \textbf{151.95} / CPI 0.966 & VAC $-$5.13 & 이내(예외 0) \\
품질 & 0.4건/FP & 1,187건 \ra{} \textbf{0.094건/FP} & 기준의 24\% & 이내 \\
리스크 & 잔여 EMV 8 / 단일 2 & 최대 잔여 7.26(9월) & — & 이내 \\
\end{tbl}
\begin{itemize}
\item 편익 — 첫 측정 B-08 13.6h, 판정은 G5 \ra{} ⑳ 원장으로 이관 · 지속가능성 — 사건 1건(P1-02) 기록
\item 9월 as-is 전망 06-11(+10) \ra{} EX-01 대안 1로 회복 — 0영업일은 회복의 결과다
\end{itemize}
\keymsg{품질을 결함 수로 적지 마라 — 분모 12,585·분자(통합\til{}UAT)·기준 0.4가 있어야 1,187이 읽힌다}
\note{교재 4.2 표. 여섯 열(기준선/실적/편차/허용범위/판정/참고 BL-01) 중 슬라이드는 다섯을 보인다. 범위 +185 FP는 BL-01 대비이며 허용 ±372의 49.7\% — 재확인 트리거 ±3\% 미발동. 품질 총 발견 2,563(설계 210·스프린트 1,166·통합\til{}UAT 1,187)이고 밀도의 분자는 통합시험\til{}UAT만이다 — 9월 전망 0.09 적중. 리스크 여유 최소 0.74(11월). 워크북 §3.3 항목 2. 예상 소요 8분.}
\end{frame}

\begin{frame}{원가 3층 — "16.05 미사용"과 "VAC $-$5.13"이 동시에 참}
\fig[0.60]{../그림/PM_Day4/fig_4_1_three_layers_final.pdf}
\figcap{그림 4-1. 승인 168.00 / 집행 한도 156.00 / PMB 146.82 대 최종 AC 151.95. EAC 이력 9호 150.86 \ra{} 11호 152.40(경고) \ra{} 17호 151.95, 예외 진입 0회}
\keymsg{PM의 자는 PMB다 — 종료 보고서는 세 층을 분리해 적고 "절감"이라는 말을 쓰지 않는다}
\note{교재 4.2 그림 4-1. 승인 예산에는 스폰서만 쓰는 관리예비비 12.0이 들어 있어 처음부터 PM의 돈이 아니고, 한도와 PMB 사이의 컨틴전시는 리스크가 실현되면 쓰라고 둔 돈이다. 세 줄 — 승인 대비 16.05 미사용(관리예비비 12.0 + 컨틴전시 반납 4.05), 한도 대비 4.05 여유, PMB 대비 VAC $-$5.13. 강조할 것: 9월 조정 CPI 0.968이 최종 0.966을 0.002 안에서 예측했다 — Day3에서 착시 1(라이선스 4.70)을 걷어 냈기 때문이며, 걷어 내지 않은 CPI 1.032로 EAC1 141.23을 채택했다면 최종은 "예상 못 한 10.7억 초과"로 기록됐을 것이다. 워크북 §3.3 항목 2·4. 예상 소요 8분.}
\end{frame}

\begin{frame}{변경·예비비 총괄 — CR-01\til{}14와 컨틴전시 10.20의 분해}
\begin{columns}[T]
\begin{column}{0.46\textwidth}
\subhead{변경 로그 총괄}
\begin{itemize}
\item 승인 10 · 조건부 1(CR-09) · 거버넌스 1(CR-10) · 기각 1(CR-05) · 이관 1(CR-12)
\item Day3 이후 — CR-12 P1 범위 외 이관 · CR-13 R-10 재정제 0.45 · CR-14 패치 0.20
\item FP 누적 +185(+1.5\%, 허용의 49.7\%), SI +1.02
\end{itemize}
\end{column}
\begin{column}{0.52\textwidth}
\subhead{컨틴전시 10.20 = 다섯 조각}
\begin{itemize}
\item 편입 1.02(변경 대가 \ra{} PMB) + 리스크 배정 집행 2.50
\item 미배정 풀 0.77 + 편차 흡수 1.86 + \textbf{반납 4.05}
\item 소진율(편입 제외) 50.3\% · 관리예비비 12.0 요청 0·\textbf{반납 12.0}
\item 유보 3건 — 클라우드 3.0 반납 / PMO 3.0 승인 전 / 전환 2.0 감액
\end{itemize}
\end{column}
\end{columns}
\keymsg{모든 집행에 결재(PM + 재경)와 CCB 월 보고가 붙어 있다 — 윤태호의 질의에 이 표가 답이다}
\note{교재 4.2 후반, 그림 4-2(보조). 리스크 배정 집행 2.50의 내역 — R-07 0.15·R-05 0.42·R-03 0.50·R-10 0.45·R-19 0.60·R-20 0.18·R-13 0.20. 미배정 풀 0.77 — 단말 0.30·아카이브 0.20·부하 시나리오 0.20·보안 0.05·임시 라이선스 0.02. 편차 흡수 1.86은 리스크 배정 없는 AC 초과(클라우드 0.76·망분리·DR 0.50·PMO 0.60)이며 이것이 VAC의 본체다. 관리예비비 대상 R-01·R-14·R-17 미실현(R-01은 분리 의결로 흡수). "컨틴전시를 합계로만 적는다"가 가장 흔한 오류 — 어느 리스크에 어느 결재로 나갔는지가 없으면 내부감사 질의에 답할 수 없다. 워크북 §3.3 항목 3·4. 예상 소요 8분.}
\end{frame}

\begin{frame}{계약 종결 — 검수 \ra{} 정산 \ra{} 유보금 \ra{} 하자보수 \ra{} 지체상금}
\begin{tbl}[\scriptsize]{L{0.14\textwidth}L{0.50\textwidth}L{0.28\textwidth}}
\textbf{단계} & \textbf{SI 한빛디지털} & \textbf{근거 서면} \\ \midrule
검수 & 최종 검수 확인서 05-26 & 검수 확인서 R1\til{}R5 \\
정산 & 69.22 + 대기 0.60 + 컨틴전시 귀속 1.20 = \textbf{71.02}, 잔금 05-31 & 합의서 7호(05-26) \\
유보금 & 5\% = 3.46억 \ra{} 하자보수보증증권 대체 & 계약 §16 · 보증증권 \\
하자보수 & 1년 2027-05-29\til{}2028-05-28, 범위 = \textbf{인수 산출물의 결함} & 개선·요구 변경은 P5 \\
지체상금 & \textbf{0 — 근거 네 줄}(다음 프레임) & 합의서 6호 · Day3 ⑮ 항목 6 \\
\end{tbl}
\begin{itemize}
\item 나머지 — 전환 용역 10.05(하자 6개월) · 상용SW 19.42 · 클라우드 22.16 중 P5 승계 · 단말 8.90
\item 하도급(모빌커넥트 220 FP) — 수행사 정산 내 1.21, \textbf{직접 지급 요청 없음} 확인
\end{itemize}
\keymsg{정산액을 한 숫자로 적지 마라 — 71.02가 69.22 + 0.60 + 1.20임을 모르면 0.60은 "왜 줬나"가 된다}
\note{교재 4.3. 순서가 곧 논리다 — 검수 없이 정산 없고, 정산 없이 유보금 없다. 하자보수 범위 정의가 없으면 하자보수 1년은 무상 유지보수 1년이 된다(결함 vs 개선의 판정은 CAB). 하도급 직접 지급 여부를 확인하지 않으면 종료 후 정산이 다시 열린다. 상용SW 잔금 4.70은 통합시험 통과 01-15에 발생 — Day3 착시 1의 해소 시점. 워크북 §3.3 항목 7. 예상 소요 8분.}
\end{frame}

\begin{frame}{지체상금 0의 네 줄 — 같은 표가 대기 비용 0.60도 설명한다}
\begin{itemize}
\item ① G4 05-28 = 기준선 완료일 \ra{} 계약 §18의 "지연" 불성립
\item ② 순 지연 분해(Day3 ⑮ 항목 6) — 발주 7 + 동시 3 + 수행 3 $-$ 회복 3 = 10 \ra{} EX-01로 0
\item ③ 수행사 단독 3일(S10 재작업) \ra{} 회복 3일(S12\til{}S17)과 상계 — 변경협의체 10-14
\item ④ 동시 3일 — 07-15 합의(SCL 원칙 10) 공기 연장 인정·비용 불인정
\item 같은 귀책 분해표의 반대 방향 = 발주사 귀책 대기 비용 0.60 지급(R-05 0.42·R-20 0.18)
\end{itemize}
\keymsg{네 줄이 없으면 0은 "봐줬다"로 읽히고, 있으면 "계산했다"로 읽힌다 — Day3 동시기록이 종료에서 정산된다}
\note{교재 4.3 그림 4-3(보조). 0은 "없음"으로 읽히기 쉬운 숫자이므로 근거가 필요하다 — 내부감사는 "왜 물리지 않았는가"를 묻는다. 핵심 문장 — 같은 표가 양쪽(지체상금 0과 대기 비용 지급)을 다 설명할 때만 그 표는 정산의 근거가 된다. 이것이 Day3 5.5절 동시기록 문화의 종료 시점 회수다. 귀책을 그때그때 적지 않았다면 종료 시점에 이 네 줄은 쓸 수 없다. 워크북 §3.3 항목 7. 예상 소요 6분.}
\end{frame}

\begin{frame}{문서·형상 인수 — 2023년 지적사항의 회수}
\begin{itemize}
\item 조건은 "받았다"가 아니라 \textbf{발주사 저장소에서 클린 빌드가 재현되는가}
\item 온담 — 재현 100\%(05-12, 12개 모듈 + 하도급 어댑터 220 FP), SBOM 312(GPL 계열 0)
\item 빌드·배포 절차서 v1.0 · 일 1회 동기화 \ra{} G4 후 발주사 저장소 단독 정본
\item 아카이브 \textsf{governance/p1-closure/} — 헌장\til{}편익 원장, 계약 서면 32건, 보존 10년
\item 지적 대응 요약 — 2020 편익 근거 \ra{} ⑳ 원장 / 2020 손실 확정 \ra{} §14 / 2023 소스 \ra{} 형상 인수
\end{itemize}
\keymsg{종료 보고서는 지적사항에 답하는 자리이지, 지적사항이 없었던 척하는 자리가 아니다}
\note{교재 4.4. 소스 파일을 받는 것은 2023년 지적("위탁 산출물 형상관리·소스 인수 절차 미비")의 상태 그대로다. 형상 인수 = SAC 항목 9의 증빙이자 헌장 §14 ②의 판정 근거. SBOM이 없으면 라이선스 리스크가 운영 중에 발견된다. 실무 오류 — 하도급분이 빠진다, 아카이브 위치가 PM의 노트북이다. 워크북 §3.3 항목 8. 예상 소요 6분.}
\end{frame}

\begin{frame}{미해결 사항 이관 목록 — 소유자 실명 없는 이관은 소멸이다}
\begin{tbl}[\scriptsize]{L{0.04\textwidth}L{0.42\textwidth}L{0.18\textwidth}L{0.14\textwidth}L{0.12\textwidth}}
\textbf{\#} & \textbf{내용} & \textbf{소유자} & \textbf{기한} & \textbf{승인} \\ \midrule
1 & 아카이브 접근 절차 재경 승인(SAC 11 조건부) & 재경팀장 & 06-15 & 승인 \\
3 & 발주 마감 22:00 운용 전환, CAB 재상정(CR-09 ①) & 한동석 & 2027-09 CAB & \textbf{승인 전} \\
6 & 이천·안성 검사기록 디지털화 — 편익 B-02 분모의 전제 & 나영선 & 09 이사회 & \textbf{승인 전} \\
7 & SAP ECC 2027-12 EoM 대응, IF-SAP 8본(KE-04) & 박세연 & 08-31 & \textbf{승인 전} \\
10 & 하자보수기 PMO 잔류 1명 예산 1.05(Day2 유보 3.0) & 프로그램 PMO장 & 06-01 운영위 & \textbf{승인 전} \\
\end{tbl}
\begin{itemize}
\item 10건 = 승인 6 · 승인 전 4 — 세 칸이 결정적: \textbf{소유자 실명 · 기한 · 승인 여부}
\item 출처 열이 앞선 열아홉 문서의 어느 행이 프로젝트 밖으로 나가는지를 기록한다
\end{itemize}
\keymsg{"승인 전"은 종료 보고서에서 가장 짧은 줄이지만 가장 오래 남는 줄이다}
\note{교재 4.5 표(10건 중 5건 발췌 — 승인 전 4건 전부 + 1번). 나머지 승인 5건: 2 미동의 9점 2차 배포(오정근 07-31) · 4 3PL 재협상(한동석 06-30) · 5 앱 벤더 계약 갱신(오정근·박세연) · 8 클라우드 사이징·유보 3.0 반납(박준영·재경팀장) · 9 ISMS-P 증적(최영주). 이관 항목은 두 종류다 — 소유자와 예산이 있어 실행만 남은 것(승인)과, 소유자는 있으나 결정이 없는 것(승인 전). 후자를 "승인"처럼 적으면 이사회 통과 순간 아무도 책임지지 않는 항목이 된다. 6번을 특히 짚는다 — 범위 밖이지만 편익의 전제이므로 적는다. 적지 않으면 G5 편익 미달이 "P1 결함"으로 기록된다. 워크북 §3.3 항목 6. 예상 소요 8분.}
\end{frame}

\begin{frame}{이관 10건의 출처와 수신 — 다른 넷과 묶는 열}
\fig[0.60]{../그림/PM_Day4/fig_4_4_followon_map.pdf}
\figcap{그림 4-4. 출처(SAC·SLA·KEDB·변경 로그·컨틴전시 유보·편익 원장·정본 제외 범위) \ra{} 항목·소유자·기한 \ra{} 수신 조직. 승인 6 실선·승인 전 4 파선, 6번 \ra{} ⑳ BR-05}
\keymsg{출처 열이 없으면 종료 시점에 누군가 떠올린 목록이고, 있으면 앞선 문서의 기록이다}
\note{교재 4.5 그림 4-4. 어디에서 오는가 — SAC 조건부(1), SLA 강등 재상정 조건(4), KEDB 의존·협의(2·7), 변경 로그 이관·조건부(3), 컨틴전시 유보(8·10), 편익 원장 전제(6), 회사 정본 제외 범위(7). 수신 열에는 결정 회의체(CAB·프로그램 이사회·운영위)가 적혀 있다 — 종료 보고서가 이사회를 통과하는 것과 승인 전 4줄이 결정되는 것은 별개다. 실무 오류 네 가지 — 소유자가 부서명, 승인 여부 열 없음, 범위 밖은 적지 않음, 출처 열 없음. 예상 소요 5분.}
\end{frame}

\begin{frame}{잔여 리스크의 운영 이관과 팀 해산 — 사람의 전환}
\begin{columns}[T]
\begin{column}{0.50\textwidth}
\subhead{잔여 리스크 — 등록부 v1.4(G4 종결)}
\begin{itemize}
\item 종결 12 / 운영 이관 3 / 편익 이관 1
\item R-13 \ra{} KE-01(오정근) · R-16 \ra{} KE-02(한동석) · R-18 \ra{} 앱 벤더(오정근·박세연)
\item R-12 \ra{} BR-01(재무본부·한동석·나영선)
\item 신규 운영 리스크 3 — 3PL 폴백 영구화·SAP EoM·IPO 변경 동결 창
\end{itemize}
\end{column}
\begin{column}{0.48\textwidth}
\subhead{팀 해산 — 사람별 일자}
\begin{itemize}
\item 발주사 14명 — 12명 06-01 복귀, \textbf{2명 P5 전환}
\item 수행사 05-31 철수, 하자보수 대기 4명 비상주
\item PMO 3명 — 2명 프로그램 PMO, 1명 잔류(승인 전, 이관 10)
\item 하자보수 창구 — P5 서비스운영팀(박준영) \ra{} UC ②, 판정 CAB
\end{itemize}
\end{column}
\end{columns}
\keymsg{리스크는 사라지지 않는다 — 운영 등록부·KEDB·편익 리스크 표 중 한 곳에 새 소유자와 함께 다시 적힌다}
\note{교재 4.6. 세 곳 어디에도 적히지 않은 잔여 리스크는 소멸이고, 실현될 때 "아무도 몰랐던 문제"로 돌아온다. 신규 운영 리스크 3건은 운영 리스크 등록부(박준영)로, 검토는 월 1회 서비스 리뷰. 2명의 P5 전환은 SAC 항목 5(교육)의 다른 얼굴 — 만든 사람이 운영 조직에 남는 것이 가장 짧은 지식 이전이다. 실무 오류 — 잔여 리스크 일괄 "종결", 해산을 한 날짜로, 하자보수 창구가 없어 첫 하자 문의가 전 PM의 개인 전화로, 잔류 인력 예산 미결정 채 잔류. Bridges의 "끝냄"은 팀에도 필요하다. 워크북 §3.3 항목 5·8. 예상 소요 7분.}
\end{frame}

\begin{frame}{문서의 마지막 독자 — 3년 뒤의 내부감사·분쟁·다음 PM}
\begin{tbl}[\scriptsize]{L{0.16\textwidth}L{0.38\textwidth}L{0.38\textwidth}}
\textbf{독자} & \textbf{묻는 것} & \textbf{답하는 곳} \\ \midrule
내부감사(3년 뒤) & 컨틴전시 R-03 0.50은 누가 언제 승인했나 & 항목 4 결재 08-12 · CCB 10-15 소급 · 11월\til{} 자동 안건 \\
분쟁(하자보수기) & 이것이 결함인가 개선인가 · 지체상금 0은 봐준 것인가 & 하자보수 범위 정의 · 검수 확인서 R1\til{}R5 · 네 줄 \\
다음 PM(P2·P5) & 온담은 컷오버를 어떻게 했는가 & 런북 실행 기록 열 · 교훈 L-15 \\
\end{tbl}
\begin{itemize}
\item 약어를 풀어 쓴다 — RB-3이 무엇인지 3년 뒤에는 아무도 모른다
\item 출처를 적는다 — 어느 문서 어느 행
\item \textbf{판단의 근거를 결론과 분리해 적는다} — "충족"이 아니라 "조건부 충족, 06-15, 재경팀장"
\end{itemize}
\keymsg{결론만 남은 문서는 3년 뒤 검증할 수 없고, 검증할 수 없는 문서는 신뢰되지 않는다}
\note{교재 4.7\til{}4.8. 06월 프로그램 이사회는 첫 독자이지 마지막 독자가 아니다. 이 장의 오류는 전부 "선언과 증명을 구별하지 못하는 것"으로 수렴한다 — 종료 시점에 처음 쓰고, 원가를 한 층으로("16억 절감"), 지체상금 0을 "없음"으로, 소스를 "받았다"로, 소유자를 부서명으로, 잔여 리스크를 일괄 종결, 해산을 날짜 하나로. 그 결과는 2019년 세 줄과 같은 종류의 문서이고, 마지막 독자는 그것을 처음부터 다시 만들어야 한다. 예상 소요 6분.}
\end{frame}

\begin{frame}{3교시 핵심 정리 — 종료 보고서는 증명이다}
\begin{itemize}
\item 8항목 · 모든 숫자에 출처 · 완료 기준 다섯에 판정자 실명 · 첫 문장 "재기준선 0회"
\item 원가 3층 분리 — 16.05 미사용 / 여유 4.05 / VAC $-$5.13 동시에 참, PM의 자는 PMB
\item 정산 71.02 = 69.22 + 0.60 + 1.20 · 지체상금 0의 네 줄 · 클린 빌드 100\%·SBOM 312
\item 이관 10건(승인 6·승인 전 4) — 소유자 실명·기한·승인 여부 없으면 소멸
\item 잔여 리스크는 새 소유자와 함께 다시 적힌다 · 해산은 사람별 일자
\end{itemize}
\keymsg{4교시 실습 ① — 워크북 §3 ⑱ 종료 보고서·이관 목록·계약 종결을 온담 값으로 직접 쓴다}
\note{제4장 핵심 정리 다섯 줄. 4교시 실습으로 넘기기 전에 워크북 §3.4 빈 템플릿과 §3.5 완성 예시본의 값이 이 교시의 숫자와 같음을 확인시킨다(교재와 워크북이 다르면 워크북이 정본). 질문 시간 포함. 예상 소요 8분.}
\end{frame}


\section{4교시 (90분) — 실습 ① 컷오버 런북·SAC + 운영 이관·SLA 3층}

% ===== 4교시 (90분) — 실습 ① 컷오버 런북·SAC + 운영 이관·SLA 3층 =====
\begin{frame}{이 교시에서 하는 것 — 오전 이론 둘을 손으로 쓴다}
\begin{itemize}
\item \textbf{⑯ 런북·SAC} — go/no-go 3 · PoNR · 롤백 트리거 4 · SAC 10항목을 \textbf{직접} 쓴다
\item \textbf{⑰ 운영 이관} — SLA/OLA/UC 3층 매핑 7행 + error budget 환전을 \textbf{손계산}한다
\item 두 문서는 서로의 입력 — SAC 항목 7(SLA 서명)이 ⑰ 항목 4, 런북 D+1 서명이 ⑰ 항목 1
\item 도구는 종이 캔버스 · 계산기 · 워크북 §1.4·§2.4 빈 템플릿 — 노트북은 검산용
\item 시간 — 안내 10분 / ⑯ 35분 / ⑰ 30분 / 짝 교환 검산 15분
\end{itemize}
\keymsg{정답은 워크북 §1.5·§2.5에 있다 — 보지 말고 쓰고, 다 쓴 뒤에 대조하라}
\note{워크북 §1·§2 실습. Day3 4교시와 같은 진행 — 과제 지시 두 프레임, 작성 요령 네 프레임, 흔한 오류 한 프레임, 예시본 대조, 정리. 오전 2교시(런북 구조·게이트·PoNR·롤백 비대칭)와 3교시(SAC·ELS·SLA 3층·error budget)가 이 90분의 재료다. 두 과제를 하나의 교시에 묶은 이유 — 런북의 마지막 행(D+1 07:20 인수 서명)이 운영 이관의 첫 행이고, SAC 항목 7(SLA/OLA/UC 서명본)은 ⑰ 항목 4의 표가 없으면 판정할 수 없다. 오픈은 이벤트가 아니라 이관이라는 부제를 두 문서가 손으로 증명하게 된다. 예상 소요 3분.}
\end{frame}

\begin{frame}{과제 ⑯ — 온담 1차 컷오버, 창 02-26 20:00\til{}02-28 06:00(34h)을 런북으로}
\begin{itemize}
\item 게이트 3개 — 착수(19:30) · 데이터 정본 전환(01:00) · 매장 배포(12:00) — 조건·판정자·재실행 허용 횟수
\item PoNR 시각과 근거, 기술적 롤백 최종 시한 — \textbf{두 시각을 분리}해 적고 그 사이의 결정권자를 적는다
\item 롤백 트리거 4 — 관측 지표 · 임계값(리허설 실측에서) · 관측 시각 · 판정자 · 전체/부분 롤백
\item SAC 10항목 — 조건(관측 가능) · 판정자(인수자 실명) · 증빙 · G3 사전 점검 · G4 판정
\item 주어진 것 — 리허설 3회 실측(이력 적재 6h 50m) · 매장 620점 · 인터페이스 47본 · PG 7본 · 헌장 §14 완료 기준 ①
\end{itemize}
\keymsg{묻는 것은 하나 — 새벽 3시에 이 종이만 보고 "지금 돌리면 되는가"에 답할 수 있는가}
\note{워크북 §1.4 빈 템플릿 항목 3·4·5·8을 쓴다. 항목 1·2(문서 정보·시각표 38행)와 항목 6·7·9·10은 예시본으로 읽기만 한다 — 시각표를 다 쓰면 35분이 끝난다. 학생이 자주 묻는 것 — "게이트는 왜 세 개인가?" 되돌리는 비용이 0 / 시간 / 시간+손실+매장으로 달라지는 지점이 셋이기 때문이다(2교시 그림 2-2). "SAC 10항목이 워크북의 12항목과 다른가?" 워크북 예시본은 12항목이고 과제는 10항목 이상이면 된다 — 운영 조직의 관심 순서(문서·모니터링·백업복구·보안·교육·KEDB·SLA·절차·소스·인터페이스·계약·잔여 결함)를 지키는지를 본다. 임계값은 반드시 숫자로 — "이상 없음"은 채점하지 않는다. 예상 소요 4분(안내), 작성 35분.}
\end{frame}

\begin{frame}{과제 ⑰ — SLA 7항목을 3층으로 받치고, 99.5\%를 분과 릴리스 횟수로 환전한다}
\begin{itemize}
\item 3층 매핑표 — 행 = SLA 7항목(가용성·P1·P2 복구·배치·발주 마감·출고 확정·데이터 정정)
\item 열 = SLA 목표 / OLA(사내 팀·목표) / UC(외부 계약·목표) / 합산(가장 긴 사슬) / 판정
\item 주어진 UC 초안 — 수행사 유지보수 착수 1h·복구 4h · CSP 99.9\% · 3PL 대성로지스 \textbf{UC 없음}
\item error budget — 99.5\%/월 \ra{} 허용 다운(분) \ra{} 비계획 60\% / 릴리스 40\% \ra{} 릴리스당 기대 다운 \ra{} 월 허용 횟수
\item 주어진 DORA 실측(2026-09) — 변경 실패율 12\% · 평균 복구 240분 / ELS 목표 10\% · 45분
\end{itemize}
\keymsg{묻는 것은 둘 — 어느 행이 구조적으로 불가능한가, 복구 시간을 얼마로 줄여야 주 1회 릴리스가 SLA와 양립하는가}
\note{워크북 §2.4 빈 템플릿 항목 3·4를 쓴다. 항목 2(운영 문서 12종)·5(ELS 8주)·6\til{}9는 예시본으로 읽는다. 3층 매핑은 7행 중 ② P1 복구와 ⑥ 출고 확정 두 행이 연습의 핵심 — 나머지 다섯은 성립하는 행이고 10분 안에 끝낸다. error budget은 계산기로 — 평균 달 30.4375일을 쓰라고 미리 말한다(30일로 계산한 조는 216분이 나오고 그것은 틀린 것이 아니라 다른 전제다). 학생이 자주 묻는 것 — "계획 정지는 다운인가?" SLA 산정 제외, 단 CAB 승인 건만 — 이 규칙 한 줄을 표에 적었는지 본다. 예상 소요 3분(안내), 작성 30분.}
\end{frame}

\begin{frame}{작성 요령 ① — PoNR은 "무엇이 되돌릴 수 없게 되는가"로 정한다}
\fig[0.58]{../그림/PM_Day4/fig_2_2_reversal_cost.pdf}
\figcap{그림 2-2. 되돌리는 비용의 계단 — 게이트 1 전 0 \ra{} 시간만 \ra{} 시간+손실(PoNR 01:40) \ra{} +매장 방문 2.1억(게이트 3) \ra{} 불가(16:00). 계단이 오르는 자리마다 결정권자가 한 층 올라간다}
\keymsg{PoNR = 신 원장이 정본이 되어 이후 거래가 그곳에만 쌓이기 시작하는 시각 — 배포 시작이 아니다}
\note{교재 §2.3, 워크북 ⑯ 항목 4. 온담은 게이트 2 통과(01:38) 직후 신 재고·주문 원장 활성화 01:40이 PoNR — 03:00 3PL 폴백 배치 1회차가 신 원장 기준으로 출고 지시를 만들기 때문에, 그 뒤의 롤백은 수작업 대사 1,200건 이상·2영업일을 남긴다. 기술적 롤백 최종 시한 16:00은 별개의 계산(오픈 06:00 $-$ 복원 8h $-$ 검증 2h $-$ 매장 캐시 재동기화 4h)이다. 두 시각을 하나로 적는 조가 매 기수 절반이다 — PoNR 이후 16:00까지는 "가능하지만 손실 있음"이고 결정권자가 P1 PM에서 정미란으로 올라간다는 것을 표에 적게 한다. 매장 배포를 PoNR로 적은 조에는 "매장은 되돌릴 수 있다, 원장 분기는 되돌리기 어렵다"를 되묻는다. 예상 소요 6분.}
\end{frame}

\begin{frame}{작성 요령 ② — SAC의 판정자는 인수자다, 프로젝트가 판정하면 자기 평가다}
\fig[0.58]{../그림/PM_Day4/fig_3_2_sac_progress.pdf}
\figcap{그림 3-2. SAC 12항목 × 2시점 — G3 사전 점검 01-28(충족 2·미해소 7·대기 3) \ra{} 14주 해소 \ra{} G4 05-21(충족 11·조건부 1). 조건부 1건은 해소 기한·책임자 서명으로 이관 1번이 된다}
\keymsg{헌장 §14 완료 기준 ① — "SAC 미해소 0건 또는 해소 기한·책임자 서명" — 이 한 줄이 판정 규칙이다}
\note{교재 §3.2, 워크북 ⑯ 항목 8. 판정자 열을 먼저 채우게 한다 — 원칙은 인수자 박준영, 기술 항목(소스·인터페이스)은 박세연, 보안은 최영주, 계약 이관은 강현도. 판정자 열에 P1 PM이나 수행사 PM이 들어간 조는 그 행을 다시 쓰게 한다. 증빙 열에 "완료 보고"라고 쓴 조도 마찬가지 — 문서 번호·시험 기록·서명본이어야 한다. G3 사전 점검 열이 있어야 G4에서 처음 보는 일이 없다는 것, 그리고 조건부 1건(항목 11 아카이브 접근 절차, 재경 승인 06-15)이 완료 기준 ①의 두 번째 절로 통과해 ⑱ 이관 1번이 된다는 것이 이 그림의 요점이다. 예상 소요 5분.}
\end{frame}

\begin{frame}{작성 요령 ③ — 3층은 아래에서 위로: UC 없는 SLA는 약속이 아니다}
\fig[0.58]{../그림/PM_Day4/fig_3_4_sla_chain.pdf}
\figcap{그림 3-4. SLA ② P1 복구 2시간 — 초안 사슬 15+30+30+240 = 315분 > 120 / 재협상 165 / 병렬 착수로 $\le$ 120. 항목 ⑥ 출고 확정은 3PL UC 부재로 "최선 노력" 강등}
\keymsg{약속은 그것을 받치는 약속의 합보다 강할 수 없다 — 합산은 가장 긴 사슬로, 한 층이 비면 "구조적 불가"}
\note{교재 §3.5, 워크북 ⑰ 항목 4. 쓰는 순서를 뒤집게 한다 — SLA 목표를 먼저 적고 OLA·UC를 끼워 맞추는 것이 아니라, UC(수행사 복구 4h, CSP 99.9\%, 3PL 없음)와 OLA(착수 30분·원인 분리 30분)를 먼저 적고 합을 올려 SLA가 성립하는지 판정한다. ② P1 복구는 초안 UC로 315분이라 구조적으로 불가 — 선택지는 UC 재협상(착수 30·복구 90, 24/7 대기, 연 +0.6억 [추가 설정])이거나 SLA 문안 강등이다. ⑥ 출고 확정 D+1 06:00은 3PL 대성로지스 UC가 물류본부 계약이라 P1 밖 — "최선 노력"으로 강등하고 UC는 "배치 파일 수신 확인 30분"으로 한정 서명(한동석 05-10). 판정 열에 "성립"만 있고 "구조적 불가"가 하나도 없는 조는 UC를 읽지 않은 것이다. 예상 소요 6분.}
\end{frame}

\begin{frame}{작성 요령 ④ — error budget 산식: 99.5\%를 분으로, 분을 릴리스 횟수로}
\begin{tbl}[\scriptsize]{L{0.30\textwidth}L{0.34\textwidth}L{0.28\textwidth}}
\textbf{단계} & \textbf{산식} & \textbf{온담 값} \\ \midrule
월 허용 다운 & 30.4375일 × 24 × 60 × 0.5\% & \textbf{219.2분} (3h 39m) \\
비계획 유보 / 릴리스 몫 & 60\% / 40\% & 131.5분 / \textbf{87.7분} \\
릴리스당 기대 다운(실측) & 변경 실패율 12\% × 복구 240분 & 28.8분 \ra{} 월 \textbf{3.0회} \\
릴리스당 기대 다운(중간) & 10\% × 90분 & 9.0분 \ra{} 월 9.7회 \\
릴리스당 기대 다운(ELS 목표) & 10\% × 45분 & 4.5분 \ra{} 월 \textbf{19.5회} \\
소진 시 규칙 & 유보 초과 시 표준 릴리스 동결 · 다음 달 초기화 & 결정자 박준영 \\
\end{tbl}
\keymsg{복구 시간을 4시간에서 45분으로 줄이지 않으면 주 1회 릴리스도 SLA와 양립하지 않는다 — 회귀 자동화·KEDB의 존재 이유}
\note{교재 §3.6, 워크북 ⑰ 항목 3. 검산값은 src/day4\_closing.py 출력과 같다. 계산 순서가 곧 전제다 — 평균 달(30.4375일)을 쓰는지, 계획 정지(2차 컷오버 8h·월 1회 유지보수 창)를 SLA 산정에서 제외하는 규칙을 적었는지, 릴리스 몫을 몇 \%로 두는지. 릴리스당 기대 다운 = 변경 실패율 × 평균 복구 시간이라는 한 줄이 CAB의 "안정 vs 기능" 감정 토론을 산수로 바꾼다. 3월 ELS 실측 57분(유보의 43\%, 예산 소진 26\%)을 표 아래 적어 예산이 "쓰이는" 것임을 보인다. 그림 3-5(도넛)는 노트판에서 보조로 쓴다. 예상 소요 6분.}
\end{frame}

\begin{frame}{흔한 오류 — 이 실습에서 매 기수 나오는 다섯}
\begin{enumerate}
\item 게이트가 하나(오픈 직전)뿐이고 조건이 "이상 없음" — 되돌리는 비용이 달라지는 지점이 게이트다
\item PoNR을 "배포 시작"으로 적고 기술적 시한과 한 시각으로 합친다 — 사이의 결정권자가 사라진다
\item 롤백 트리거가 "중대 장애 발생 시" — 임계값 없는 트리거는 새벽에 토론이 된다
\item SAC 판정자가 프로젝트 쪽, 증빙이 "완료 보고" — 인수가 아니라 자기 평가
\item SLA만 있고 OLA는 구두, 3PL·앱 벤더는 UC에서 뺀다 — 99.5\%를 분으로 환산하지 않는다
\end{enumerate}
\keymsg{다섯 오류의 공통점 — 판정자와 숫자를 적지 않았다. 판정자 없는 조건은 조건이 아니다}
\note{워크북 §1.3·§2.3 각 항목의 "흔한 오류" 줄을 모은 것. 짝 교환 검산 15분에서 이 다섯을 체크리스트로 준다 — 상대 조의 런북에서 ① 게이트 수와 조건의 관측 가능성 ② PoNR과 기술적 시한이 다른 시각인가 ③ 트리거 4개 모두에 숫자 임계값이 있는가 ④ SAC 판정자 열에 프로젝트 사람이 있는가 ⑤ 3층 표에 "구조적 불가" 행이 있는가·219.2분이 나왔는가. 세 개 이상 걸리는 조는 예시본 대조 전에 그 항목만 다시 쓴다. 예상 소요 4분.}
\end{frame}

\begin{frame}{예시본 참조 · 시간 배분 · 심화 과제}
\begin{itemize}
\item 예시본 — 워크북 §1.5(런북 38행·게이트 3·트리거 4·SAC 12) · §2.5(문서 12종·3층 7행·error budget·ELS 8주)
\item 인쇄용 — \textsf{../템플릿/PM\_Day4/16\_컷오버런북\_SAC\_완성예시.pdf} · \textsf{17\_운영이관\_SLA\_완성예시.pdf} (빈 양식은 \textsf{\_빈템플릿.pdf})
\item 시간 — ⑯ 게이트·PoNR·트리거·SAC 35분 / ⑰ 3층 매핑·환전 30분 / 짝 교환 검산·대조 15분
\item 자기 점검 — 워크북 §1.7·§2.7 점검 질문 각 3개. 값이 다르면 \textbf{어느 전제가 달랐는지}를 적는다
\item 심화(3\til{}7년차) — 자기 회사의 최근 오픈 당일 기록에서 게이트·PoNR·임계값·편차 기록 넷을 찾아 본다
\end{itemize}
\keymsg{예시본은 정답이 아니라 전제가 적힌 답이다 — 다른 값이 나왔으면 전제를 비교하라}
\note{Day3 4교시와 같은 운영. 예시본 PDF는 실습 시작 시 배포하지 않고 35분 뒤(⑰ 착수 시점)에 ⑯ 예시본만 먼저 배포, ⑰ 예시본은 손계산이 끝난 뒤. 워크북 §1.7 기본 과제(게이트 2에서 재실행 2회 후에도 14건 불일치 \ra{} RB-1 발동 \ra{} 항목 3·4·5·7 재작성, R-14 확률 갱신, 롤백 런북 10행)와 §2.7 기본 과제(3PL 재협상 타결 시 3층 표 재작성·릴리스 몫 40\% 유지 가능성 계산)는 시간이 남는 조의 추가 과제이고, 대개 과제로 낸다. 심화 과제의 네 질문 중 셋 이상이 "아니오"인 조직은 다음 오픈에서 이번 오픈으로부터 아무것도 배우지 못한 채 시작한다는 문장을 그대로 읽어 준다. 예상 소요 90분 전체.}
\end{frame}

\begin{frame}{4교시 핵심 정리}
\begin{enumerate}
\item 게이트 3개는 되돌리는 비용이 달라지는 세 지점 — 조건은 관측값, 판정자는 실명, 재실행 횟수까지
\item PoNR 01:40과 기술적 시한 16:00은 다른 시각 — 그 사이는 "가능하지만 손실 있음", 결정권자가 올라간다
\item 트리거 4개의 임계값은 리허설 실측에서 온다 — 발동 0이 결과이지 목표가 아니다
\item SAC는 인수자가 판정해야 인수다 — 조건부 1건은 기한·책임자 서명으로 이관 목록에 넘어간다
\item 3층은 아래에서 위로, 99.5\%는 219.2분 \ra{} 87.7분 \ra{} 릴리스 3회 — 복구 시간이 릴리스 속도를 정한다
\end{enumerate}
\keymsg{5교시는 종료 — 보고서는 선언이 아니라 증명이고, 미해결 사항은 소유자 실명이 없으면 소멸한다}
\note{조별 산출물 회수. 3층 매핑표 1장은 6교시 시작 전까지 벽에 붙여 두고 갤러리워크로 서로 대조하게 한다 — "구조적 불가" 행의 수가 조마다 다르면 그것이 5교시 미해결 사항 이관(어느 약속이 프로젝트 밖으로 넘어가는가)의 도입 재료다. 예상 소요 3분.}
\end{frame}

\section{5교시 (75분) — 교훈과 편익: 프로젝트가 끝난 뒤에 남는 것}

% =====================================================================
% 5교시 — 교훈과 편익 (교재 제5장, 워크북 §4~§5)
% =====================================================================
\begin{frame}{5교시에서 배우는 것 — 프로젝트가 사라진 뒤의 독자를 위해 쓴다}
\begin{itemize}
\item 교훈이 조직에 남지 않는 이유 — 수집이 아니라 \textbf{흡수}의 실패 (Williams · Syllk)
\item 교훈 등록부 6칸 — 적용 대상·소유자 없는 교훈은 감상 / 온담 15건 \ra{} 개정 요청 3건
\item 편익의 책무 분리(ITO) — output은 G4에서 종결, outcome은 2029-05까지 실명 서명
\item 분모 정치 91.3 vs 84.6 — 승인된 목표는 승인된 분모로 판정한다
\item 손익분기 66\%·두 시계·편익 리스크 5건 — 그리고 P1은 어느 성공인가
\end{itemize}
\keymsg{교훈과 편익의 공통 실패는 하나 — 프로젝트가 사라진 뒤의 독자를 위해 쓰지 않은 것이다}
\note{제5장 도입. 4교시까지가 "종료를 증명하는 문서"였다면 5교시는 "종료 뒤에도 살아 있어야 하는 문서" 둘 — ⑲ 교훈 등록부와 ⑳ 편익 실현 원장 — 을 다룬다. 둘의 독자는 P1 PM이 아니라 06-01 이후의 프로그램 PMO장·재경팀장·오정근 본부장이며, 그래서 모든 칸에 PM이 아닌 실명이 들어간다. 데이터 일자 2027-05-28(G4), 원장 v1.0 서명 05-21, 개정 요청 발송 06-04. 워크북 §4·§5 완성 예시본이 정본이고 교재와 다르면 워크북을 따른다. 예상 소요 4분.}
\end{frame}

\begin{frame}[shrink=8]{교훈은 왜 조직에 남지 않는가 — Williams(2008)·Syllk의 여섯 층}
\fig[0.6]{../그림/PM_Day4/fig_5_1_syllk.pdf}
\figcap{그림 5-1. Syllk 6층(사람 3 / 시스템 3)과 교훈 6칸. ④ 적용 대상·⑤ 소유자가 사람 층에서 시스템 층으로 가는 문이고, 개정 요청서가 그 문을 통과하는 유일한 형식 — 6개월 소멸선은 [추가 설정]}
\keymsg{교훈이 남는 조건은 교훈의 수가 아니라 — 그 교훈으로 개정된 양식이 다음 프로젝트의 착수 패키지에 있는가다}
\note{교재 5.1절. Williams(2008, PMJ 39(4))의 실증 — 교훈 학습의 실패는 모으지 못해서가 아니라 모은 교훈이 절차(양식·규정·계약 문안)로 흡수되지 않아서다. Duffield \& Whitty(2015, IJPM)의 Syllk 모형은 흡수 조건을 여섯 층으로 구조화한다 — 사람 쪽 learning·culture·social, 시스템 쪽 technology·process·infrastructure. 회고 워크숍과 PM 커뮤니티는 사람 층만 통과시키고, 시스템 층에 닿지 않으면 6개월 뒤 사라진다. 이것은 Day1 1.3절 임시조직 이론과 같은 말이다 — 임시조직이 해산하면 배운 것도 해산하므로 해산 전에 영구조직의 제도로 옮겨야 한다. 온담이 15건마다 적용 대상을 세 표준 개정 요청으로 묶는 이유가 여기 있고, 개정 요청서는 등록부의 부속이 아니라 결과물이다. 틀리는 방식 셋 — 저장소에 넣고 끝, "검토 바람", 회고를 의식으로. 예상 소요 6분.}
\end{frame}

\begin{frame}{교훈 등록부의 여섯 칸 — "적용 대상"이 없는 교훈은 감상이다}
\begin{tbl}[\scriptsize]{L{0.14\textwidth}L{0.42\textwidth}L{0.36\textwidth}}
\textbf{칸} & \textbf{무엇을 적는가} & \textbf{규율} \\ \midrule
① 상황 & 무엇이 있었나 — 사실·일자·숫자 & 형용사 없이 \\
② 원인 & 왜 그랬나 — 양식·규칙·구조 & \textbf{사람 이름 금지} \\
③ 교훈 & 무엇을 바꾸면 되나 — 한 문장 & "중요하다"는 감상 \\
④ 적용 대상 & 어느 양식·규칙·문안의 어느 칸 (NP / STD / CTR) & "전사"는 없음과 같다 \\
⑤ 소유자 & 누가 바꾸나 — 실명 & PM은 06-01에 없어진다 \\
⑥ 증빙 & 어느 문서의 어느 행 & 워크북 ⑫\til{}⑯ 항목 번호 \\
\end{tbl}
\begin{itemize}
\item 잘된 것도 교훈이다 — 우연한 성공을 양식으로 고정하지 않으면 다음 프로젝트가 다시 발명해야 한다
\end{itemize}
\keymsg{"담당자가 놓쳤다"는 사실이지 교훈이 아니다 — "양식에 그 칸이 없었다"로 적어야 칸이 다음 담당자를 잡는다}
\note{교재 5.2절, 워크북 §4 항목 1. 여섯 칸이 다 채워질 때만 교훈이고 하나라도 비면 감상이다. 원인 칸에 사람을 적지 않는 것이 이 설계의 규율 — 다음 프로젝트의 담당자도 놓칠 것이므로, 양식이 바뀌어야 그 뒤의 모든 프로젝트에서 칸이 잡는다. 사람 이름은 소유자 칸에만 나온다. 적용 대상 약호 NP 다음 프로젝트 양식 / STD 사내 표준 / CTR 계약 문안. 온담 Lessons log 41건 중 종료 정리 15건(나머지 26: CIR 12·P3/P4 소관 9·재발 낮음 5), 원천 분포 조기 신호 5 / 착시 1 / 감리·계약 3 / 이슈·리스크 2 / 절차·거버넌스 3 / 잘됨 1(L-15; L-08도 잘됨). 검산 15 = 5 + 1 + 3 + 2 + 3 + 1. 예상 소요 6분.}
\end{frame}

\begin{frame}{온담 P1 교훈 15건 중 다섯 — 원인에 사람이 없고 소유자에 PM이 없다}
\begin{tbl}[\scriptsize]{L{0.07\textwidth}L{0.24\textwidth}L{0.30\textwidth}L{0.17\textwidth}L{0.12\textwidth}}
\textbf{ID} & \textbf{상황} & \textbf{원인(구조) \ra{} 교훈} & \textbf{적용 대상} & \textbf{소유자} \\ \midrule
L-02 & 승인 대기 12영업일, 대기 비용 0.42 & 헌장·RACI에 대체 결정권자 칸 없음 \ra{} 실명 지정 & STD ① 헌장 §11 & PMO장 \\
L-06 & 감리 AF-1-03·04 귀속 다툼, 0.18 & 감리 계약에 귀속 판정 기준 없음 \ra{} WBS 담당 열·RACI & CTR 감리 부속서 3 & 강현도·정하늘 \\
L-08 잘됨 & 대기 비용 판정 1회로 종결 & 회의록 "요청·응답 일자" 열 \ra{} 동시 기록 서명 & CTR §14 / NP 회의록 & 강현도·P1 PM \\
L-09 & D$-$2 라이선스 CFO 결재 대기, 0.02 & 창 앞 사전 결재 규정 없음 \ra{} 4주·5천만 원 일괄 결재 & STD ② 내부통제 & 재경팀장·정미란 \\
L-15 잘됨 & 두 컷오버 정시·트리거 0 & 1차 지표를 2차 go 조건으로 \ra{} 컷오버 표준 & STD ① 컷오버 표준 & PMO장·박준영 \\
\end{tbl}
\keymsg{L-02의 원인을 "본부장이 바빴다"로 적으면 사실이지만 교훈이 아니다 — 다음 프로젝트의 Senior User도 바쁠 것이다}
\note{교재 5.3절, 나머지 10건은 워크북 §4.5. 다섯 건은 Day3·Day4의 사건이 계약 문안과 사내 표준으로 바뀌는 경로를 보여 준다 — L-02는 Day3 ⑬ IS-02·⑫ 신호 2, L-06은 Day3 4.3절 "등급보다 귀속 칸이 무겁다"가 계약 부속서가 되는 길, L-08은 동시 기록(contemporaneous records) 덕에 4.3절의 지체상금 0이 네 줄로 증명된 근거, L-09는 2.8절 D$-$2 사건, L-15는 R-14 트리거를 게이트 조건으로 묶은 사슬 — 어느 문서에도 "표준"으로 적혀 있지 않았다. L-08의 소유자에 P1 PM이 있는 것은 06-01에 프로그램 PMO장으로 승계된다는 뜻이다. 실무에서 틀리는 방식 다섯 — 원인에 사람, 적용 대상 "전사", 소유자 PM, 잘된 것 없음, 개정 요청이 부속. 예상 소요 7분.}
\end{frame}

\begin{frame}{사내 표준 개정 요청 3건(06-04 발송) — 15건을 15곳에 보내면 15건이 "검토 중"이 된다}
\begin{tbl}[\scriptsize]{L{0.22\textwidth}L{0.17\textwidth}L{0.12\textwidth}L{0.41\textwidth}}
\textbf{표준} & \textbf{소유자} & \textbf{의결 시점} & \textbf{개정 조항(교훈)} \\ \midrule
① 프로젝트 관리 규정·양식집 & 프로그램 PMO장 & 2027-07 운영위 & 성과보고 SV·CV 분해 열(L-01·03) · 헌장 대체 결정권자(L-02·10) · RoC 부속(L-11) · 컷오버 표준(L-15) · BC 편익 표 분모·실명 필수(⑳) \\
② 재무본부 내부통제 지침 & 재경팀장·정미란 & 2027-08 & 5,000만 원 조항 단서 신설 · 개인정보 사건 판단 시한 절차(L-09·14) \\
③ 구매 계약 표준 문안 & 강현도·최영주 & 2027-09 P2 후속 계약부터 & SI 계약 §14 여유 소유권·동시 지연·대기 기록 · 감리 부속서 3 귀속 기준 · 조달 상태 양식(L-06\til{}09) \\
\end{tbl}
\begin{itemize}
\item 조항 번호와 새 문장이 있는 요청만 의결된다 — "검토 바람"은 영원히 검토된다
\item 적용 시점이 P2 G2 전 — 다음 프로젝트의 착수 패키지에 들어가야 흡수다
\end{itemize}
\keymsg{개정 요청서는 교훈 등록부의 부속이 아니라 결과물이다 — 아카이브에 함께 들어가면 함께 잊힌다}
\note{교재 5.3절 후반, 워크북 §4 항목 2. 세 표준은 각각 소유 부서와 개정 절차가 하나이므로 15건을 세 묶음으로 보내면 세 번의 의결로 끝나고, 15곳에 보내면 15건이 "검토 중"이 된다. 각 요청서에는 현행 조항 번호·새 문장·근거 교훈 ID·적용 시점이 있고, ①은 P2 G2 전, ③은 P2 후속 계약부터 적용된다. 지속적 성공(Cooke-Davies)의 판정 기준이 바로 이 세 개정이 07·08·09월에 실제 의결되어 P2·P5·2028 후속 과제 착수 패키지에 들어갔는가이다 — 마지막 프레임에서 회수한다. 예상 소요 5분.}
\end{frame}

\begin{frame}[shrink=14]{편익 실현의 책무 분리 — Zwikael \& Smyrk의 ITO, 측정 책임자 실명이 서명이다}
\fig[0.58]{../그림/PM_Day4/fig_5_2_ito_ledger.pdf}
\figcap{그림 5-2. Input \ra{} Transform \ra{} Output $|$ Utilisation \ra{} Outcome과 책임 분리선. output 서명(P1 PM·이재현)은 G4 인수로 종결, outcome(오정근 5·한동석 1)·utilisation 서명은 2029-05까지 유효(BR-04) — 원장 6칸의 ⑤ 칸이 서명란}
\keymsg{P1 PM은 폐기율을 내릴 수 없다 — PM이 outcome에 서명하면 없어진 조직이 책임을 지고, 그것은 아무도 지지 않는 것과 같다}
\note{교재 5.4절, 워크북 §5 항목 5. Day1 3장의 ITO 모형 회수 — output(플랫폼)의 책임은 프로젝트에, outcome(폐기율 2.2\%)의 책임은 프로젝트 소유자(현업 책임자, PRINCE2 Senior User)에, 그 사이 utilisation(매장이 발주 가이드를 지키고 구매팀이 폐기를 집계하는 것)이 있다. 이 분리가 없으면 G5에서 편익 미달의 원인은 "P1이 잘못 만들었다"와 "현업이 안 썼다" 사이에서 영원히 다투어진다. ITO 서명표 05-21: output P1 PM·이재현 / outcome 오정근(B-04·05·07·03·08)·한동석(B-02) / utilisation 매장운영팀장·강현도·마케팅·인사본부·실사 담당·나영선·P3 PM·박준영·재경팀장 / 판정 재무본부(정미란·재경팀장). 서명자 변경 시 후임이 재서명한다. Day1 BC §4 "측정 책임자를 부서명이 아니라 실명으로"가 이 날을 위한 것이었다. 반대 실패 — outcome 소유자가 서명하지 않으면 "IT 시스템의 성과"가 되어 현업은 쓰지 않아도 면책된다. 예상 소요 6분.}
\end{frame}

\begin{frame}[shrink=14]{편익 실현 원장 6칸 — 직접 편익 118 / 온담 귀속 99.6억}
\begin{tbl}[\scriptsize]{L{0.07\textwidth}L{0.30\textwidth}L{0.24\textwidth}L{0.15\textwidth}L{0.13\textwidth}}
\textbf{B-ID} & \textbf{편익(측정식·리포트)} & \textbf{기준값 \ra{} 목표} & \textbf{측정 책임자} & \textbf{승인 / 귀속} \\ \midrule
B-04 & 폐기율 = 월 폐기 금액 $\div$ 취급원가 (RPT-INV-05) & 3.9\%(2025) \ra{} 2.2\%(M+24) & \textbf{오정근} / 강현도 & 43 / 43 \\
B-05 & 배달수수료 회피 = 전환 매출 $\times$ 12.8\% (RPT-ORD-02) & 앱 2.1\% \ra{} 온라인 18\% & \textbf{오정근} / 재경팀장 & 37 / \textbf{18.6} \\
B-07 & 직영 인건비율(분기) & 21.4\% \ra{} 20.1\% & \textbf{오정근} / 인사본부 & 38 / 38 \\
B-02 & 재고 정확도, 두 분모 병기 (StRS-034) & A 91.3 \ra{} 98.5 / B 84.6 미설정 & \textbf{한동석} / 실사 담당 & (B-04 기여) \\
B-08 & 가맹 발주 리드타임 (RPT-FRN-03) & 38h \ra{} 12h · \textbf{첫 측정 13.6h} & \textbf{오정근} / 박준영 & (B-03 기여) \\
\end{tbl}
\begin{itemize}
\item 6칸 = 측정식(코드 수준) · 기준값(값·시점·분모·출처) · 목표값 · 측정 시점 · 측정 책임자 실명 · 미달 시 조치
\item 미달 시 조치는 "원인 분석"이 아니라 output / outcome / utilisation 세 갈래 중 하나로 분류 \ra{} 그 갈래의 소유자가 움직인다
\end{itemize}
\keymsg{원장 v1.0 2027-05-21 서명, 정본 관리자는 재무본부 — P1 종료 후에도 살아 있어야 하는 유일한 P1 문서다}
\note{교재 5.4절 표, 워크북 §5 항목 1. B-03 결품률(4.7 \ra{} 1.5\%)은 기회손실 산정식이 2027-11 확정이라 금액 미산정으로 표에서 뺐다. 검산 118 = 43 + 37 + 38, 온담 귀속 99.6 = 43 + 18.6 + 38 — B-05의 37은 이사회 승인본 숫자이지만 가맹 매장의 배달 매출은 가맹점주의 것이므로(C-13은 연결매출이 아니다) 온담 손익 귀속은 직영 비중 2,980/5,920 = 50.3\%인 18.6억, 나머지 18.4억은 가맹점주 편익이다. B-08 진행률 94\% = (38 $-$ 13.6)/(38 $-$ 12). 미달 시 조치의 분기를 B-04로 예시 — output(리포트 오류) \ra{} 하자보수 / outcome(발주 정확도·수요 예측 미사용) \ra{} 오정근의 발주 가이드 채택 조치 / utilisation(매장 미준수) \ra{} P3 교육·지역매니저 점검. 미달 판정 시 오정근이 G5 전 개선 계획을 낸다. 예상 소요 6분.}
\end{frame}

\begin{frame}[shrink=14]{분모 정치 — 91.3 vs 84.6, 승인된 목표는 승인된 분모로 판정한다}
\fig[0.58]{../그림/PM_Day4/fig_5_3_denominator.pdf}
\figcap{그림 5-3. A 순환실사 SKU(69\%) 91.3 \ra{} 98.5(7.2\%p) vs B 전체 SKU 84.6(13.9\%p, 1.93배). 비순환 31\%의 검사기록은 종이·엑셀(2025-11 이사회 반려 34억) — 판정 A 서명 05-21, B 목표는 이관 6번 가결/부결(BR-05) 뒤 운영 +12에서 설정}
\keymsg{Day3의 정정(R-12)은 실은 편익 판정권의 문제였다 — 누가 어느 분모로 98.5를 판정하는가를 서명으로 확정하지 않으면 G5는 협상이 된다}
\note{교재 5.6절, 워크북 §5 항목 2. 분모 A = 순환실사 SKU 기준 91.3\%(물류본부 2025 보고, O-08), 분모 B = 전체 SKU 기준 84.6\%(2026-12-18 StRS-034 두 분모 리포트 첫 실행 — Day2에서 Must로 둔 이유). 차이의 원인은 비순환 SKU(원료·반제품·포장재)의 입고 검수·검사기록이 종이·엑셀이라 실사 정확도 자체가 낮고, P1은 그 디지털화를 하지 않았다(헌장 §5 제외 범위). 논리 셋 — ① 이사회 승인본의 분모는 "순환실사 SKU"로 명기 ② B로 바꾸면 개선 폭 13.9\%p, 약속하지 않은 것을 약속한 셈 ③ B의 정확도는 P1 밖의 전제 위에 있다. 합의 한동석·나영선·재경팀장 05-21, 이사회 06월 보고에 "두 분모 병기·판정은 A" 명시. 유리한 분모를 택한 것이 아니라 약속한 분모로 판정하고, 약속하지 않은 분모의 전제(이관 6번, 나영선 2027-09 부의)를 정직하게 적은 것이다. 틀리는 방식 넷 — 분모를 안 적음, 유리한 분모 사후 선택, 불리한 분모로 "정직" 과시(전제 미충족이 P1 결함으로 기록됨), 전제 미이관. 예상 소요 7분.}
\end{frame}

\begin{frame}{편익이 실패하는 이유 — Ward \& Daniel의 셋과 온담의 세 가지 정치}
\begin{columns}[T]
\begin{column}{0.48\textwidth}
\subhead{실패 원인 셋 (Ward \& Daniel 2012)}
\begin{itemize}
\item 소유자 부재 — "회사"의 편익은 아무의 것도 아니다
\item 측정 시점 부재 — M+24가 달력의 어느 날인가
\item 기준값의 정치 — 사후에 정해진 기준값은 협상이 된다
\end{itemize}
\end{column}
\begin{column}{0.50\textwidth}
\subhead{온담 원장의 세 정치}
\begin{itemize}
\item 분모 정치 — A 91.3 / B 84.6, 판정은 A
\item 귀속 정치 — B-05 37 \ra{} 온담 18.6 + 가맹점주 18.4
\item 측정 조직 소멸 — B-08 로그: P1이 만들고 P5가 읽고 영업본부가 해석
\end{itemize}
\end{column}
\end{columns}
\keymsg{기준값은 프로젝트 전에, 측정 책임자는 프로젝트 중에, 판정은 프로젝트 후에 — 셋이 같은 문서에 있어야 한다}
\note{교재 5.5절, 워크북 §5.1. 편익 의존성 네트워크 — IT 산출물과 편익 사이에 반드시 업무 변화가 있고 그 변화의 소유자가 편익의 소유자다(ITO의 utilisation을 다른 방향에서 본 것). 귀속 정치를 적어 두지 않으면 G5에서 재무본부가 손익을 대조하는 순간 "37억이 어디 있는가"가 되고 그때 처음 설명하면 P1 편익 전체가 의심받는다. 적어 두면 18.4억은 협의회(서정필)에 "가맹점주 편익"으로 공유할 수 있는 숫자가 되어 부속서 데이터 용도 논쟁(R-02·BR-03)의 지렛대가 된다. 측정 조직 소멸은 세 조직 중 하나가 바뀌면 측정이 끊기는 리스크 — BR-04로 로그 보존 3년·서명자 변경 시 재서명·정본 재무본부. 틀리는 방식 — 승인 금액 전액 귀속, "M+24"만 적고 달력 날짜·시계 없음, 로그 보존 정책 없음. 예상 소요 5분.}
\end{frame}

\begin{frame}[shrink=11]{손익분기 66\% · 62.1\% · 78\% — 그리고 M+12 / M+24 리뷰의 두 시계}
\fig[0.56]{../그림/PM_Day4/fig_5_4_breakeven_clocks.pdf}
\figcap{그림 5-4. NPV 대 실현률(50\% $-$58.3 / 66\% 0 / 70\% +14.8 / 100\% +124.4; AC 151.95 반영 시 손익분기 62.1\%·+138.3), 판정선 95\%·66\%, 온담 귀속 기준 78\% + 프로그램 시계(M+24 2028-01·G5 2028-03-31)와 운영 시계(+12 2028-05·+24 2029-05)}
\keymsg{BC의 시계는 고치지 않는다 — 승인 문서를 사후에 고치는 것은 기준값을 고치는 것과 같다. 대신 운영 시계를 추가한다}
\note{교재 5.7절, 워크북 §5 항목 3·4. 손익분기 66.0\% = 118 $\times$ 0.66 = 77.8억/년(할인율 8\%, 운영비 27억 전액 P1 귀속, 보수적). 최종 AC 151.95 반영 시 62.1\% — 컨틴전시 반납 4.05·관리예비비 12.0 미사용의 효과가 3.9\%p 하락이다. 온담 귀속 99.6 기준으로는 요구 실현률이 78\%(77.8/99.6)로 올라간다 — 이것도 원장에 적혀야 G5가 정직하다. 판정 규칙(허용범위 $-$5\%p): 95\% 미만 경고 / 66\% 미만은 헌장 §14 중단 조건 수준 — 프로젝트는 종료되었으므로 프로그램 이사회의 후속 투자 판단 근거. 두 시계 — BC는 프로그램 착수 2026-01 기준이라 B-08의 M+12가 컷오버 전(2027-01)이 되는 모순이 있지만 고치지 않고, 운영 이관 2027-05 기준 시계를 추가한다(운영 +12 원장 v2: B-02 분모 B 목표 설정·측정 조직 확인, +24 원장 v3: 유지 판정·원장 종결 또는 상시 KPI 전환). 첫 측정 2027-04 B-08 / 06 B-04·02·03 / 09 B-05·07 / 11 B-03 산정식. 예상 소요 6분.}
\end{frame}

\begin{frame}{편익 리스크 5건 — 원장의 부속, 자금원은 운영 예산·프로그램 예비비}
\begin{tbl}[\scriptsize]{L{0.08\textwidth}L{0.26\textwidth}L{0.06\textwidth}L{0.48\textwidth}}
\textbf{ID} & \textbf{리스크} & \textbf{P} & \textbf{대응} \\ \midrule
BR-01 & 분모 분쟁 재점화(R-12 이관) & 30\% & 05-21 합의문·이사회 보고에 "판정은 A" 명시 \\
BR-02 & 외생변수 — 원재료가·이상 기후로 폐기율 왜곡 & 40\% & 수량 기준 병기, 원재료가 지수 통제 \\
BR-03 & B-05 귀속 분쟁 — 협의회가 18.4억을 본사 편익으로 오해 & 25\% & 귀속 분리 명기·가맹 귀속 지표 공유 \\
BR-04 & 측정 조직 소멸 — 로그·읽는 조직·해석 조직 셋 중 하나 변경 & 35\% & 로그 보존 3년·서명자 변경 시 재서명·정본 재무본부 \\
BR-05 & 전제 미충족 — 검사기록 디지털화 부결 시 B 목표 불가 & \textbf{50\%} & 이관 6번 부의안에 B-02·B-04 연결 명시 \\
\end{tbl}
\begin{itemize}
\item Day1 프리모템 14번 "측정 책임자가 경기 탓이라 했다" \ra{} BR-02·BR-04로 대응된 것
\end{itemize}
\keymsg{편익 리스크는 프로젝트 리스크와 다르다 — 프로젝트는 끝났고, 리스크의 소유자는 원장에 서명한 사람들이다}
\note{교재 5.7절 후반, 워크북 §5 항목 6. 편익 리스크 5건은 원장 부속이며 자금원이 프로젝트 예비비가 아니라 운영 예산·프로그램 예비비다 — P1 예비비는 G4에서 반납되었다(4.05). BR-05의 확률이 가장 높은 이유는 2025-11 이사회가 검사기록 디지털화 34억 SI 견적을 이미 한 번 반려했기 때문이고, 부결 시 B 목표는 "A와 동일 방법으로 측정 가능한 범위"로 한정하며 B-04의 원료 폐기 부분도 미실현으로 남는다. 귀인(BR-01·03)·외생(BR-02)·측정 조직 소멸(BR-04)·전제(BR-05)의 네 유형은 여러분 회사의 편익 원장에도 그대로 적용된다. 예상 소요 4분.}
\end{frame}

\begin{frame}{2027-05-28의 P1은 어느 성공인가 — Cooke-Davies의 셋과 Flyvbjerg의 곱셈}
\begin{tbl}[\scriptsize]{L{0.2\textwidth}L{0.15\textwidth}L{0.55\textwidth}}
\textbf{성공(Cooke-Davies 2002)} & \textbf{판정} & \textbf{근거} \\ \midrule
프로젝트 관리 성공 & \textbf{허용범위 이내} & 재기준선 0 · 컷오버 2회 정시 · VAC $-$5.13(한도 내, PMB 대비 3.5\% 초과) · FP 12,585 100\% · 밀도 0.094 \\
프로젝트 성공(편익) & \textbf{판정 불가} & 첫 측정 하나(B-08 13.6h, 94\%)뿐 — 판정은 G5 2028-03-31 \\
지속적 성공 & 더 뒤 & 세 표준 개정(07·08·09월)이 P2·P5 착수 패키지에 들어갔는가 \\
\end{tbl}
\begin{itemize}
\item "온버짓"이라는 말을 세 층(승인·한도·PMB) 중 어느 것으로 쓰느냐에 따라 답이 다르다 — 이 교재는 그 말을 쓰지 않는다
\item Flyvbjerg의 곱셈 — 원가·일정 두 인수는 확정, 편익 인수는 아직 곱해지지 않았다
\end{itemize}
\keymsg{종료 보고서의 편익 축 "미판정 — ⑳ 원장으로 이관"이 정직한 문장이다 — 원장·서명·두 시계·리스크는 장치이지 결과가 아니다}
\note{교재 5.8절. Day1 1.2절의 두 개념 회수 — 세 성공은 서로 독립이고, Flyvbjerg(2014)의 곱셈은 그 독립성의 통계적 표현이다. 2019년 스마트오더가 Day1 그림 1-6에서 "온타임·온버짓 + 편익 없음" 칸에 놓였던 것을 기억시켜라 — P1이 그 칸에 놓이지 않으리라는 보장은 2027-05-28 시점에 없다. 종료 보고서에 "편익 달성 예정"을 적는 것이 가장 흔한 실패이며, 그러면 G5는 판정이 아니라 협상이 되고 2020년 지적이 2028년에 되풀이된다. 5.9절 실무 실패의 공통점 — 교훈에서도 편익에서도 프로젝트가 사라진 뒤의 독자를 위해 쓰지 않은 것. 생각해 볼 질문 3·5번(여러분이 종료한 프로젝트의 편익에 여러분이 서명했는가 / P1을 "성공"이라 부를 수 있는가)을 짝과 2분 토론. 예상 소요 6분.}
\end{frame}

\begin{frame}{5교시 핵심 정리 — 교훈은 양식으로, 편익은 실명으로 남는다}
\begin{itemize}
\item 교훈 실패는 흡수의 실패 — 시스템 층에 닿는 문은 적용 대상·소유자 칸, 형식은 개정 요청서
\item 6칸이 다 차야 교훈, 원인에 사람이 있으면 교훈이 아니다 — 15건 \ra{} 3표준(07·08·09월)
\item 편익 책무는 output(G4 종결) / outcome(오정근·한동석) / utilisation — 실명 칸이 서명
\item 승인된 목표는 승인된 분모로 — A 91.3 \ra{} 98.5 판정, B 84.6은 전제 결정 뒤 운영 +12
\item 손익분기 66.0 / 62.1 / 78\% · 두 시계 · BR 5건 — P1은 관리 성공은 이루고 편익 성공은 미판정
\end{itemize}
\keymsg{프로젝트는 06-01에 없어진다 — 그 뒤에 남는 것은 개정된 양식 세 개와 서명된 원장 하나다}
\note{제5장 핵심 정리 다섯 줄. 6교시 실습 ②(워크북 §4~§5)로 연결 — ⑲ 교훈 등록부 15건 중 조기 신호 5건을 여섯 칸으로 채우고 개정 요청서 ① 조항 하나를 문장으로 쓰는 것, ⑳ 원장에서 B-02 분모 합의문과 B-05 귀속 분리를 적는 것이 과제다. 검산을 미리 알려라 — 15 = 5 + 1 + 3 + 2 + 3 + 1, 118 = 43 + 37 + 38, 99.6 = 43 + 18.6 + 38, 78\% = 77.8/99.6, B-08 진행률 94\%. §4.7 기본 과제(G2 중단 시 premature closure의 교훈 10건)는 1.2절과 이 장의 결합, §5.7 기본 과제(G5 첫 12개월 실적으로 미달 시 조치 실행)는 5.4\til{}5.7절의 결합 문제. 예상 소요 3분.}
\end{frame}

\section{6교시 (75분) — 실습 ② 종료 보고서·미해결 이관·교훈·편익 원장}

% =====================================================================
% 6교시 — 실습 ②
% =====================================================================
\begin{frame}{6교시 — 이 교시에서 하는 것}
\begin{itemize}
\item 5교시의 세 문서(⑱·⑲·⑳)를 \textbf{손으로} 쓴다 — 워크북 §3\til{}5
\item 세 과제, 한 실 — 이관 6번(나영선) \ra{} 교훈 적용 대상 \ra{} 원장 BR-05
\item ⑱ 종료 보고서 최종 4축 + 이관 목록 5건 이상 (25분)
\item ⑲ 교훈 5건 — 여섯 칸, 적용 대상·증빙까지 (20분)
\item ⑳ 편익 원장 B-01\til{}B-06 × 6칸, 측정 책임자 실명 (25분)
\end{itemize}
\keymsg{종료는 선언이 아니라 증명이다 — 세 문서의 모든 숫자·이름은 앞 문서의 어느 행에서 왔는지 추적된다}
\note{Day3 6교시와 같은 운영 — 개인 작성 후 조별 대조, 강사는 순회. 다른 점은 세 과제가 "프로젝트가 끝난 뒤에도 남는 것"을 적는다는 것이다 — 이관 목록은 P1이 하지 않은 것, 교훈은 다음 프로젝트가 바꿀 것, 원장은 운영이 잴 것. 세 문서는 한 실로 이어진다: ⑱ 이관 6번(입고 검사기록 디지털화, 나영선, 2027-09 이사회)이 ⑲ 개정 요청 ①(BC 편익 표에 분모 칸)의 근거가 되고 ⑳ B-02 분모 B의 목표 설정 조건(BR-05)으로 다시 나타난다. 수강생에게 워크북 §3.7·§4.7·§5.7 기본 과제와 §3.5·§4.5·§5.5 완성 예시본을 펴게 한다. 시간이 부족하면 ⑲는 3건, ⑳은 B-02·B-04·B-07 세 항목으로 줄이되 6칸은 다 채운다. 예상 소요 3분.}
\end{frame}

\begin{frame}{과제 ⑱ — 워크북 §3: 종료 보고서 최종 4축 + 이관 목록}
\begin{itemize}
\item 시점 2027-05-28 G4 — 워크북 §3.7 기본 과제 상황(SAC 항목 7 미충족)
\item 최종 4축 — 범위 FP 12,585 / 일정 G4 = 기준선 완료일 / 원가 AC 151.95 / 품질 SAC 11+1
\item 원가는 3층을 다 적는다 — 168 승인 / 156 한도 / 146.82 PMB \ra{} 여유 4.05 · VAC $-$5.13
\item 이관 목록 5건 이상 — 출처·소유자 실명·기한·승인/승인 전·수신 조직
\item 헌장 §14 ① 두 경로 중 택일 — "미해소 0" / "해소 기한·책임자 서명"
\end{itemize}
\keymsg{원가 한 층만 적으면 "절감"이 되고, 세 층을 적으면 "16.05 미사용·여유 4.05·초과 5.13"이 동시에 참이 된다}
\note{워크북 §3.3 항목 1·2·6과 §3.5 예시본 항목 2·6을 쓴다. 기본 과제는 SAC 항목 7(SLA/OLA/UC)이 UC ② 재협상 결렬로 미충족인 상황 — 06-18(허용범위 끝)까지 G4를 미루는 대안과 05-28 조건부 종료의 비용(PMO·수행사 잔류·하자보수 개시 지연·유보금)을 비교하고 항목 1·6·7·8을 다시 쓴다. 최종 4축은 그림 4-1(3층 + EAC 이력 9호 150.86 \ra{} 11호 152.40 \ra{} 17호 151.95)의 값을 그대로 인용한다 — 종료 시점에 새로 만든 숫자는 없어야 한다. 이관 목록은 예시본 10건(승인 6·승인 전 4) 중 5건 이상을 자기 문장으로 다시 쓰되, 미충족 항목 7이 "승인"인지 "승인 전"인지를 반드시 판정하게 한다(UC 예산 증액 결정자가 누구인가). 예상 소요 25분.}
\end{frame}

\begin{frame}[shrink=8]{미해결 사항 이관 10건 — 출처에서 수신 조직까지}
\fig[0.6]{../그림/PM_Day4/fig_4_4_followon_map.pdf}
\figcap{그림 4-4. 출처(SAC 조건부·CR 조건·SLA 강등·편익 전제) \ra{} 항목·소유자 실명·기한 \ra{} 수신 조직. 승인 6 실선·승인 전 4 파선. 6번(나영선·2027-09 이사회)이 ⑳ BR-05로 이어진다}
\keymsg{이관 목록에 "이 프로젝트의 편익이 실현되려면 이 프로젝트가 하지 않은 것"이 한 줄이라도 있는가}
\note{워크북 §3.5 항목 6 표의 10건을 지도로 옮긴 것. 승인 전 4건 — 3번(가맹 마감 22:00 전환, 한동석, 2027-09 CAB), 6번(입고 검사기록 디지털화, 나영선, 2027-09 이사회), 7번(SAP ECC 유지보수 종료 대응, 박세연, 별도 과제), 10번(PMO 잔류 예산, 프로그램 PMO장, 06-01 운영위) — 은 무엇이 승인되어야 하는지(누가·어느 회의체)를 적었기 때문에 종료와 함께 사라지지 않는다. 특히 6번은 P1 범위 밖(종이·엑셀 검사기록)이지만 재고 정확도 98.5\% 목표의 전제이므로 적었다 — 2025-11 이사회가 34억 SI 견적을 반려한 이력까지. 수강생의 목록에서 소유자가 부서명이거나 기한이 "추후"인 행을 찾아 표시하게 한다. 예상 소요 5분.}
\end{frame}

\begin{frame}{과제 ⑲ — 워크북 §4: 교훈 5건, 여섯 칸}
\begin{itemize}
\item 상황 · 원인 · 교훈 · \textbf{적용 대상} · 소유자 · \textbf{증빙} — 여섯 칸을 다 채운다
\item 상황은 숫자로 — 12영업일 · 0.42억 · +38분 · D$-$2 발견
\item 원인은 사람이 아니라 양식·규칙·구조 — "그 칸이 없었다"
\item 5건 중 잘된 것 1건 이상(L-08 동시 기록 · L-15 분리 컷오버 연결)
\item 5건을 개정 요청 3건(규정·내부통제 지침·계약 문안)으로 묶는다
\end{itemize}
\keymsg{"적용 대상"이 없는 교훈은 감상이다 — 문서명과 칸까지 적히고 그 문서의 소유자가 접수해야 교훈이다}
\note{워크북 §4.3 항목 1·2와 §4.5 예시본(15건)을 쓴다. 기본 과제는 헌장 §14 중단 조건이 2026-11-27 G2에서 발동했다고 가정한 premature closure의 교훈 10건인데, 시간상 이 교시에서는 정상 종료 교훈 5건으로 줄여 진행하고 중단 시나리오는 과제로 남긴다. 교재 §5.3의 다섯 건(L-02 대체 결정권자, L-06 감리 귀속 기준, L-08 동시 기록, L-09 컷오버 창 앞 사전 결재, L-15 분리 컷오버의 연결)을 본으로 삼되, 수강생은 자기 문장으로 다시 쓴다. 채점의 핵심은 원인 칸에 사람 이름이 있는지, 소유자가 PM인지(PM은 곧 없어진다), 증빙이 워크북 산출물의 항목·행까지 적혔는지 세 가지다. 잘된 것의 교훈이 없으면 다음 프로젝트는 이번 성공을 다시 발명해야 한다. 예상 소요 20분.}
\end{frame}

\begin{frame}{교훈이 사라지는 경로 — Syllk 6층과 여섯 칸}
\fig[0.6]{../그림/PM_Day4/fig_5_1_syllk.pdf}
\figcap{그림 5-1. 사람 3층 / 시스템 3층, 6개월 소멸선. ④ 적용 대상·⑤ 소유자 칸이 process·infrastructure로 가는 문이다 — 개정 요청 3건(07·08·09월)이 그 문을 통과한다}
\keymsg{교훈은 아카이브에서 죽는다 — 개정 요청서는 등록부의 부속이 아니라 결과물이다}
\note{교재 §5.1 Williams(2008)와 Duffield \& Whitty의 Syllk. 사람 층(learning·culture·social)에 머문 교훈은 팀 해산과 함께 6개월 안에 소멸한다는 것이 그림의 소멸선이다. 시스템 층(technology·process·infrastructure)으로 넘어가는 유일한 문이 여섯 칸의 ④ 적용 대상과 ⑤ 소유자다. 온담 P1은 06-04에 개정 요청 3건을 발송했다 — ① 프로젝트 관리 규정·양식집(프로그램 PMO장, 2027-07 운영위) ② 재무본부 내부통제 지침(재경팀장 정미란, 컷오버 창 앞 사전 결재 목록) ③ 구매 계약 표준 문안(강현도·최영주, 감리 귀속·동시 지연·대기 기록). 수강생의 5건이 이 셋 중 어디로 가는지 표시하게 한다. 예상 소요 4분.}
\end{frame}

\begin{frame}{과제 ⑳ — 워크북 §5: 편익 원장 B-01\til{}B-06 × 6칸}
\begin{itemize}
\item 6칸 — 측정식 · 기준값 · 목표값 · 측정 시점 · \textbf{측정 책임자 실명} · 미달 시 조치
\item 측정식은 코드 수준 — 분자·분모·산출 시스템(RPT-FRN-03)·주기
\item 기준값에는 반드시 분모와 측정 시점 — 91.3\% (순환실사 SKU, 05-21)
\item 미달 시 조치는 원인 분기 — output / outcome / utilisation
\item 직접 편익 118억, 온담 귀속 99.6억 — 두 기준 실현률을 따로 적는다
\end{itemize}
\keymsg{프로젝트가 끝난 뒤 이 원장을 들고 있는 조직과 사람의 이름이 적혀 있는가 — 그것이 편익 원장의 첫 조건이다}
\note{워크북 §5.3 항목 1·2·5와 §5.5 예시본을 쓴다. 온담 P1의 6항목 — B-04 폐기율 43억·B-05 배달수수료 37억·B-07 인건비율 38억·B-02 재고 정확도·B-03 결품률·B-08 리드타임(참고 2 — B-01 온라인 비중은 B-05로 환산, B-06 처리량은 P2 이관). 과제명은 B-01\til{}B-06 여섯 칸 표이나 항목 번호는 예시본의 B-ID를 그대로 쓰게 한다. 측정 책임자는 편익 소유자 실명 + 측정 실무 직책 — 오정근 5·한동석 1·재무본부 판정, ITO 서명표(항목 5)의 세 열과 같은 사람이어야 한다. 미달 시 조치는 "원인 분석"이 아니라 output이면 하자보수·유지보수, outcome이면 업무 변화 소유자, utilisation이면 P3 채택 조치로 분기한다. 시간이 없으면 B-02·B-04·B-07 세 항목만 6칸으로. 예상 소요 25분.}
\end{frame}

\begin{frame}[shrink=16]{ITO 사슬과 책무 분리 — ⑤ 칸은 서명란이다}
\fig[0.6]{../그림/PM_Day4/fig_5_2_ito_ledger.pdf}
\figcap{그림 5-2. Input \ra{} Transform \ra{} Output ǀ Utilisation \ra{} Outcome. 서명자 3열 — output P1 PM·이재현 / utilisation 매장운영팀장·강현도·나영선·박준영 / outcome 오정근 5·한동석 1·재무본부 판정 — 과 원장 6칸(B-04 행)}
\keymsg{output 소유자(PM)는 종료 후 편익 판정에서 빠진다 — 편익을 잰 사람이 편익을 만든 사람이면 판정이 아니다}
\note{교재 §5.4 Zwikael \& Smyrk의 ITO 모델. 책임 분리선의 왼쪽(Input·Transform·Output)은 프로젝트, 오른쪽(Utilisation·Outcome)은 운영이다. 원장 ⑤ 칸(측정 책임자)은 outcome 열의 사람이어야 하며 PM이 들어가면 안 된다 — PM은 G4 다음 날 없고, 있더라도 자기 output을 스스로 판정하게 된다. 수강생의 원장에서 측정 책임자가 "영업본부"이거나 PM이면 오류로 표시한다. utilisation 열이 비어 있는 항목(채택이 안 되면 편익이 안 나는 항목 — B-05 앱 전환)은 P3 채택 조치와 연결해야 한다. 예상 소요 4분.}
\end{frame}

\begin{frame}{작성 요령 ①②③ — 이관 한 줄 · 적용 대상 · 기준값 확정}
\begin{tbl}[\scriptsize]{L{2.2cm}L{4.6cm}L{4.6cm}}
요령 & 규칙 & 온담 P1의 한 줄 \\ \midrule
① 이관 한 줄의 5요소 & 출처 · 소유자 실명 · 기한 · 승인/승인 전(누가·어느 회의체) · 수신 조직 & 6번: 편익 원장 B-02 전제 · 나영선 · 2027-09 이사회 · 승인 전 · 생산본부 \\
② 교훈의 적용 대상 & 세 층 중 하나 이상 — 다음 프로젝트 양식 / 사내 표준 조항 / 계약 문안 · 문서명·칸까지 & L-09: 내부통제 지침 "컷오버 창 앞 4주 사전 결재 목록" 조항 신설 · 재경팀장 \\
③ 기준값 확정 & 승인본이 전제한 분모로 판정 · 다른 분모는 병기 · \textbf{이유를 적는다} & A 91.3\%(순환실사 SKU, 승인본 전제) 확정 · B 84.6\% 병기, 목표 M+12 설정 \\
\end{tbl}
\keymsg{91.3인가 84.6인가 — 어느 쪽이든 값 옆에 "왜 이 분모인가"를 적지 않으면 G5에서 판정권 다툼이 된다}
\note{워크북 §3.3 항목 6, §4.3 항목 1, §5.3 항목 2의 "어떻게 적는가"를 한 표로 압축했다. ①은 다섯 요소 중 하나라도 비면 그 행은 종료와 함께 소멸한다 — 특히 "승인 전"은 무엇이 승인되어야 하는지를 적어야 다음 회의체 안건이 된다. ②는 "전사"·"향후 프로젝트"가 적용 대상이 아니라는 뜻이다 — 개정될 문서의 이름과 조항·칸, 그 문서의 소유 부서를 적는다. ③은 Day3 §5.6 정정(91.3)이 편익 판정권의 문제였음을 다시 확인하는 자리다 — 승인본이 약속한 분모로 판정하는 것이 정직하고, 확장 분모의 목표는 P1이 하지 않은 전제(검사기록 디지털화, 이관 6번) 위에 세울 수 없다. 유리한 분모로 바꾸거나 둘 중 하나를 지우는 것이 흔한 오류. 예상 소요 6분.}
\end{frame}

\begin{frame}{분모 정치 — A 91.3 \ra{} 98.5 와 B 84.6, 판정과 목표는 다른 문서다}
\fig[0.6]{../그림/PM_Day4/fig_5_3_denominator.pdf}
\figcap{그림 5-3. A 91.3\ra{}98.5(7.2\%p) vs B 84.6(13.9\%p 필요, 1.93배). 비순환 SKU 31\%(이천·안성 원료·반제품, 검사기록 34억 반려). 판정 A 서명 05-21 \ra{} B 목표 조건(이관 6번) \ra{} 가결/부결(BR-05) 분기}
\keymsg{두 분모를 다 적고 판정 분모를 명시한다 — 하나를 지우는 순간 나머지 하나가 정치가 된다}
\note{교재 §5.6과 워크북 ⑳ 항목 2. 분모 B로 판정하면 목표까지 13.9\%p가 필요해 A의 1.93배 — 같은 프로젝트가 성공에서 실패로 바뀐다. 그래서 온담은 승인본 전제인 A로 판정 기준값을 확정하고(05-21 서명), B는 병기하되 목표는 M+12 리뷰(2028-05)에서 검사기록 과제의 이사회 결정(2027-09, 이관 6번)에 따라 설정한다 — 가결이면 B 목표 설정, 부결이면 B는 참고 지표로 남긴다(원장 BR-05). 수강생 과제 ⑳의 B-02 행에 이 분기가 적혔는지 확인한다. §5.7 기본 과제(G5 실적 B 88.1\%)는 이 분기를 실제로 실행하는 연습이다. 예상 소요 4분.}
\end{frame}

\begin{frame}{흔한 오류 5 · 예시본과 템플릿 · 시간 배분}
\begin{columns}[T]
\begin{column}{0.52\textwidth}
\subhead{흔한 오류 5}
\begin{itemize}
\item 원가 실적을 한 층(168 승인)과만 비교 — "절감 16.05"
\item 이관 소유자가 부서명, 기한이 "추후", 승인 전을 승인처럼
\item 교훈의 원인이 사람, 적용 대상이 "전사", 소유자가 PM
\item 기준값에 분모·측정 시점 없음, 유리한 분모로 교체
\item 측정 책임자가 "영업본부" · 미달 시 조치가 "원인 분석"
\end{itemize}
\end{column}
\begin{column}{0.46\textwidth}
\subhead{예시본·템플릿}
\begin{itemize}
\item ⑱ §3.4 템플릿 · §3.5 예시본 · §3.6 해설
\item ⑲ §4.4 템플릿 · §4.5 예시본 15건
\item ⑳ §5.4 템플릿 · §5.5 예시본 · 그림 5-2·5-3
\end{itemize}
\subhead{시간 배분 · 심화}
\begin{itemize}
\item ⑱ 25 / ⑲ 20 / ⑳ 25 / 대조 5
\item 심화(3\til{}7년차) — 자기 회사 종료 보고서·교훈·편익 표를 §3.7·§4.7·§5.7 심화 과제로 센다
\end{itemize}
\end{column}
\end{columns}
\keymsg{셋 중 둘 이상이 없으면 그 보고서는 종료를 선언한 문서이지 종료를 증명한 문서가 아니다}
\note{순회하며 다섯 오류를 표시하게 한다 — 대부분의 오류는 칸이 비어 있는 것이 아니라 칸이 이름 대신 부서명·"추후"·"전사"·"원인 분석"으로 채워진 것이다. 예시본은 베끼는 것이 아니라 대조하는 것 — 수강생의 값이 예시본과 다르면 워크북이 정본이지만, 다른 이유를 적을 수 있으면 그것이 더 좋은 답이다. 심화 과제 셋은 회사에 돌아가 하는 것: ① 원가가 어느 층과 비교됐는가·다른 층으로 바꾸면 절감이 초과로 바뀌는가 ② 이관 목록에 실명·승인 전 표시가 몇 개인가 ③ 지체상금 정산에 근거 문장이 있는가 / 교훈 3건의 여섯 칸 중 적용 대상·소유자를 채울 수 있는가·실제 개정된 양식이 하나라도 있는가 / 승인 문서 편익 표에서 종료 후 실제 측정된 항목이 몇 개인가(0이면 2020 감사 지적이 그 회사에도 해당). 예상 소요 5분.}
\end{frame}

\begin{frame}{6교시 핵심 정리 — 프로젝트가 끝난 뒤에 남는 세 문서}
\begin{itemize}
\item ⑱ 종료 보고서 — 모든 숫자는 성과보고·변경 로그·등록부의 행에서 온다
\item 이관 목록 — 5요소가 다 있어야 종료와 함께 소멸하지 않는다
\item ⑲ 교훈 — 상황은 숫자, 원인은 구조, 적용 대상은 문서명과 칸
\item ⑳ 원장 — 6칸, 측정 책임자는 outcome 열의 실명, 판정 분모를 명시
\item 세 문서는 한 실 — 이관 6번 \ra{} 개정 요청 ① \ra{} BR-05, 끊기면 편익은 승인 시점에만 존재한 숫자
\end{itemize}
\keymsg{끝났다고 선언한 사람·받았다고 서명한 사람·승인한 사람이 세 사람인가 — 그리고 3년 뒤 이 원장을 들고 있을 사람의 이름이 있는가}
\note{정리 3분. 세 점검 질문으로 닫는다 — ⑱: 종료 보고서의 모든 숫자가 앞 문서의 어느 행에서 왔는지 추적되는가(추적되지 않는 숫자는 종료 시점에 만들어진 숫자다) / ⑲: 원인 칸에 사람 이름이 없고 적용 대상의 문서 소유자가 개정 요청을 접수했는가 / ⑳: 편익 항목마다 output·outcome·utilisation 소유자가 세 사람이고 output 소유자가 판정에서 빠져 있는가. 마무리 30분(제6장, 그림 6-1 20종 사슬)으로 넘어가며 ⑳에서 ①로 되돌아가는 여섯 가닥(측정 책임자·분모·손익분기·AC·지체상금 0·PoNR)을 예고한다. 예상 소요 3분.}
\end{frame}

\section{마무리 (30분) — PM 트랙을 마치며 · Product 트랙 예고}

% =====================================================================
% 마무리 — 교재 제6장 · 워크북 §7
% =====================================================================
\begin{frame}{이 교시에서 — 30분}
\begin{itemize}
\item 학술 배경 6갈래 — 표준이 "하라"고만 말하는 곳의 "왜"와 "언제 틀리는가"
\item 읽을 자료 — 필수 4 / 심화 8 / 참고 / 읽지 않아도 되는 것
\item 4일 20종 산출물의 사슬을 거꾸로 — ⑳에서 ①로 (그림 6-1)
\item 세 표준의 전환기 — PMBOK 8 · ITIL v5 · 개인정보 보호법 제21445호
\item Day1\til{}4 네 문장 · Product 트랙 예고 · 사후 읽기와 30일 과제
\end{itemize}
\keymsg{오늘의 한 문장 — 오픈은 이벤트가 아니라 이관이며, 이관은 받는 쪽이 서명해야 끝난다}
\note{교재 제6장과 "PM 트랙을 마치며", 워크북 §7. 6교시 실습 ②(이관·종료·교훈·원장)를 막 끝낸 상태이므로, 이 교시는 새 내용을 쌓지 않고 4일을 하나의 사슬로 묶는다. 학술 배경은 문헌 소개가 아니라 각 문헌의 "흔한 오독" 한 줄에 집중한다 — 오독을 말해야 실무자가 기억한다. 후반 15분은 Product 트랙 예고와 30일 과제에 쓴다. 예상 소요 2분.}
\end{frame}

\begin{frame}{학술 배경 I — 책임: 누가 편익에 서명하는가}
\begin{tbl}[\scriptsize]{L{2.6cm}L{4.2cm}L{4.6cm}}
문헌 & 이 교재가 빌린 것 & 흔한 오독 \\ \midrule
Zwikael \& Smyrk (2011) & ITO 모형 · output/outcome 책임 분리 · 프로젝트 소유자 \ra{} ⑳ ITO 서명표 & "PM은 편익에 무관심해도 된다" — 아니다, utilisation이 가능한 output을 설계할 책임(두 분모 리포트 Must) \\
Ward \& Daniel (2012) & 편익 의존성 네트워크 · 편익 소유자 = 업무 변화 소유자 \ra{} 5.5절 세 실패 원인 & 편익 관리 = 종료 후 측정 — 아니다, 착수 때 네트워크를 그려 업무 변화 없는 편익을 미리 드러내는 것 \\
Williams (2008) · Duffield \& Whitty (2015) & 교훈은 수집되나 적용되지 않는다는 실증 · Syllk 6층 \ra{} ⑲ "적용 대상·소유자" 칸 & Syllk = 지식관리 도구 도입 논거 — 아니다, technology는 6층 중 하나, process·infrastructure 없으면 도구는 저장소 \\
\end{tbl}
\keymsg{outcome 소유자가 없는 프로젝트는 시작하면 안 된다 — 저자들의 실제 주장이다}
\note{교재 6.1\til{}6.3. 세 문헌은 모두 "책임"에 관한 것이다. Zwikael \& Smyrk에서 빌린 규칙은 "output 서명은 G4에서 종결된다"이고, 그래서 원장의 ⑤ 칸이 서명란이다(그림 5-2). Ward \& Daniel의 초점이 착수 시점이라는 점은 Day1 BC §4의 "측정 책임자" 열과 연결해 말한다 — 그 열이 부서명이었다면 오늘 서명할 사람이 없었다. Williams의 게재지는 리딩리스트 표기와 달라 [검증 필요]로 남겨 두었음을 정직하게 말한다. 예상 소요 4분.}
\end{frame}

\begin{frame}{학술 배경 II — 기억과 속도: 왜 교훈은 남지 않고, 안정은 환전되는가}
\begin{tbl}[\scriptsize]{L{2.6cm}L{4.2cm}L{4.6cm}}
문헌 & 이 교재가 빌린 것 & 흔한 오독 \\ \midrule
Lundin \& Söderholm (1995) & 임시조직의 Transition · 제도화된 종료 \ra{} 종료 직전 이탈과 교훈 실패의 구조적 설명 & 종료를 "행사"로 읽는 것 — 처방은 종료의 제도화, 즉 오늘의 다섯 문서 \\
Beyer 외, SRE (2016) & error budget = "안정과 속도의 환율" · SLI/SLO/SLA 구분 \ra{} 219.2분 · 릴리스 몫 87.7 & "다운을 허용하는 면죄부" — 아니다, 릴리스 동결 규칙(policy)을 숫자로 갖는 것 \\
Ford, Ford \& D'Amelio (2008) & 저항은 변화 설계에 관한 정보, 때로 추진자가 만든 것 \ra{} 3.9절 ADKAR 장벽 격자 & 저항 = 극복 대상 — 아니다, 가맹 K/A · 물류 A/D 장벽은 설계의 피드백 \\
\end{tbl}
\keymsg{함께 읽을 것 — Accelerate(변경 실패율·복구 시간)는 환전의 입력, ADKAR·Kotter·Bridges는 저항의 진단 도구}
\note{교재 6.4\til{}6.6. Lundin \& Söderholm은 Day1 1.3절에서 네 구조 개념을 빌렸던 논문이고, 오늘 Transition을 회수한다 — 4일이 한 논문의 앞뒤로 닫히는 셈이다. SRE는 3.6절의 계산 원전이며, 3장·4장은 온라인 공개라 필수 자료에 넣었다. Ford 외는 3.9절의 태도를 정한 논문이다 — 마감 22:00 저항(R)을 "가맹점주가 틀렸다"로 적지 않고 이관 3번(승인 전)으로 적은 이유가 여기 있다. 예상 소요 4분.}
\end{frame}

\begin{frame}{읽을 자료 — 필수(과정 전) · 심화(과정 후) · 읽지 않아도 되는 것}
\begin{tbl}[\scriptsize]{L{1.4cm}L{3.9cm}L{4.7cm}L{1.2cm}}
구분 & 자료 & 왜 · 어느 부분 & 시간 \\ \midrule
필수 & Day3 워크북 완성 예시본 5종 & Day4 산출물의 입력 — EX-01 대안 3 · R-14 · 귀책 분해표 · 착시 5 & 60분 \\
필수 & PRINCE2 7 Closing a Project & 5활동 · End Project Report · Follow-on action & 40분 \\
필수 & SRE ch.3 · ch.4 (온라인 공개) & 99.5\% \ra{} 분 \ra{} 릴리스 환전의 원전 & 40분 \\
필수 & 개인정보 보호법 제21445호 조문 & '유출등' 정의 · 72시간 — 직접 읽고 온다 & 30분 \\
심화 & Zwikael \& Smyrk · Ward \& Daniel · Duffield \& Whitty · Williams · Ford 외 · Humble \& Farley ch.10 · Accelerate · ITIL Foundation v5 & 6.1\til{}6.7절의 원전 — 회사 규정 개정 전에 & 25시간 \\
\end{tbl}
\keymsg{읽지 않아도 되는 것 — "종료 체크리스트 100가지"(항목 수가 아니라 세 칸이 문제다) · ITIL 판본 비교 블로그 · 편익 관리 SW 백서}
\note{교재 6.9. 필수 4종은 과정 전 자료이므로 이미 읽었어야 한다 — 읽지 않은 수강생에게는 사후 30일 과제와 함께 다시 안내한다. 심화 8종은 합계 약 25시간이며 순서를 정해 준다: 회사 규정 개정을 맡은 사람은 ITIL Foundation v5부터, 편익 책임을 맡은 사람은 Zwikael \& Smyrk부터, PMO는 Duffield \& Whitty부터. 참고 자료(ISO/IEC 20000-1 조항 · MSP 5판 Benefits · 회사 정본 §5.5·§6.5 · 검산 스크립트 day4\_closing.py)는 워크북을 쓰다 손을 뻗어 확인하는 것이라 화면에 올리지 않았다. 예상 소요 2분.}
\end{frame}

\begin{frame}[shrink=8]{4일 20종 산출물의 사슬을 거꾸로 — ⑳에서 ①로}
\fig[0.60]{../그림/PM_Day4/fig_6_1_twenty_chain.pdf}
\figcap{그림 6-1. 편익 원장 6칸 \ra{} 편익 책무 분리(ITO) \ra{} 헌장 §3 성공 기준 \ra{} BC §4 측정 책임자·분모·§6 손익분기 66\%. 여섯 가닥(측정 책임자·분모·손익분기·AC·지체상금 0·PoNR) — 고리가 비면 그 자리에서 숫자가 끊긴다}
\keymsg{20종은 스무 개의 문서가 아니라 하나의 사슬이다 — Day1에 측정 책임자를 부서명으로 적었다면 오늘 ITO 서명표에는 서명할 사람이 없었다}
\note{교재 "PM 트랙을 마치며"·워크북 §7. 그림의 여섯 가닥을 하나씩 손으로 짚는다. ⑳ 원장 "측정 책임자 실명" ← ① BC §4 열 / "분모" ← BC §4 산출 분모 + 가정 A3 / "손익분기 66\%" ← BC §6 / ITO outcome 소유자 ← ② 헌장 §3 + ④ 거버넌스(BC §9). ⑱ 최종 AC 151.95 ← ⑨ PV 곡선(145.80 \ra{} BL-02 146.82) + ⑫ EAC4 150.86 / 지체상금 0 ← ⑧ 여유 소유권 문안 + ⑮ 귀책 분해표 / 이관 6번 ← ⑦ StRS-034 두 분모 리포트 + ① A3. ⑯ PoNR·트리거 ← ⑧ 컷오버 창 허용범위 0 + ⑩ R-14 + ② 헌장 §7. 반사실 세 개를 말한다 — Day2에 두 분모 리포트를 Must로 두지 않았다면 분모 정치는 G5에서 처음 터졌고, Day3에 라이선스 4.70 착시를 걷어 내지 않았다면 최종 AC는 "예상 못 한 10.7억 초과"로 기록되었을 것이다. 예상 소요 5분.}
\end{frame}

\begin{frame}{세 표준의 전환기를 지나는 실무자에게 — 2026년 9월}
\begin{tbl}[\scriptsize]{L{3.0cm}L{4.4cm}L{4.0cm}}
표준 & 무엇이 바뀌었나 & 이 교재의 규율 \\ \midrule
PMBOK 8판 (PMI, 2025) & Focus Area · 성과영역 — 6판 번호와 병기 & 판본을 정직하게 말한다(각 X.2 대조표) \\
ITIL Version 5 (2026-02) & ITIL Value System · Lifecycle 8활동 — practice 명칭 유지 & 어휘를 왜 고정하는지 고지한다("ITIL 4가 최신" 금지) \\
개인정보 보호법 제21445호 (시행 2026-09-11) & '유출등'에 위조·변조·훼손 추가 · 72시간 통지·신고 & 조문은 단정, 적용은 비단정, 시한은 단정 \\
\end{tbl}
\subhead{회사 규정을 개정할 때의 순서}
\begin{itemize}
\item ① 지금 계약서에 적힌 어휘 확인 \ra{} ② 바뀐 상위 구조 매핑 \ra{} ③ 법령의 시한을 절차에 박는다
\end{itemize}
\keymsg{어휘가 낡았다고 규정을 통째로 바꾸는 것과, 판본이 바뀐 줄 모르는 것은 같은 실수의 두 얼굴이다}
\note{교재 "세 표준의 전환기" · 6.7 · 6.8. 이번 주에 개정법이 시행되었다는 사실을 다시 짚는다 — 3.8절의 72시간 타임라인(그림 3-6)은 그 조문 위에 있고, 인쇄 전 조문 원문 대조와 법무 검토는 이 교재가 대신하지 못한다. ITIL v5 사실은 웹 확인한 것이며 8활동 매핑표는 리딩리스트 §1.12에 [검증 필요]로 있다. 규정 개정 순서 셋은 그대로 회사에 가져가도 되는 절차다 — 어휘부터 확인하는 이유는 계약서(SLA·UC)에 적힌 어휘를 바꾸면 계약 변경이 되기 때문이다. 예상 소요 3분.}
\end{frame}

\begin{frame}{Day1\til{}4 — 네 문장}
\begin{tbl}[\small]{L{1.4cm}L{4.2cm}L{5.4cm}}
Day & 한 문장 & 오늘 어디에 남았나 \\ \midrule
Day1 & 숫자 없는 헌장은 소설이다 & 헌장 §3 성공 기준 \ra{} ⑳ ITO 서명표 · BC §6 66\% \\
Day2 & 기준선은 협상된다 & ⑧ 여유 소유권 · ⑨ 145.80 \ra{} 146.82 · ⑩ R-14 \\
Day3 & 편차는 결정 요청으로 끝난다 & ⑫ EX-01 스폰서 결정 \ra{} G4 = 기준선 완료일 \ra{} 지체상금 0 \\
Day4 & 오픈은 이관이다 & ⑯ D+1 07:20 서명 · ⑰ SLA 3층 · ⑱ 이관 10건 · ⑳ 두 시계 \\
\end{tbl}
\keymsg{네 문장은 하나다 — 결정되지 않은 것을 결정하고, 분해해 숫자로 만들고, 편차를 결정으로 되돌리고, 받는 쪽이 서명하게 한다}
\note{Day1\til{}Day3 마무리 교시의 "오늘의 한 문장"을 세로로 세운다. 각 문장이 오늘 어느 산출물에 남았는지를 말하는 것이 이 프레임의 목적이다 — 문장 자체는 이미 들었다. Day3의 문장은 "통제는 편차를 재는 일이 아니라 누가 무엇을 결정해야 하는지를 적는 일"이었고, 오늘 그 결정(EX-01 대안 1, 10-16)이 G4 = 기준선 완료일이 되어 지체상금 0의 첫 줄이 되었다(그림 4-3). 예상 소요 3분.}
\end{frame}

\begin{frame}{Product 트랙 예고 — "당신은 발주자였다. 다음 4일은 공급자다"}
\begin{columns}[T]
\begin{column}{0.48\textwidth}
\subhead{4일 동안 — 발주사 온담의 자리}
\begin{itemize}
\item "대기와 추가는 변경대가 대상"에 WBS 담당 열·RACI·동시 기록으로 답했다
\item 검수 보류 회신 · 감리 귀속 이의 · 지체상금 0과 대기 비용 0.60을 같은 표로
\item 서명한 SAC·SLA·원장 = 인도받은 것을 조직이 받아 쓰기 위한 문서
\end{itemize}
\end{column}
\begin{column}{0.50\textwidth}
\subhead{Day5\til{}8 — 반대편에 선다}
\begin{itemize}
\item 온담은 같은 회사로 다시 등장 — 이번에는 당신의 \textbf{고객}
\item 당신은 한 스타트업의 프로덕트 책임자 — 온담을 \textbf{밖에서} 본다
\item 종료 보고서를 읽을 수 없고, 이관 목록에 무엇이 있는지 모른다
\item 세 본부장의 이음새 · CFO 결재 절차 · 부채비율 187\%를 공급자 관점으로
\end{itemize}
\end{column}
\end{columns}
\keymsg{발주자로 배운 것 — 편익의 분모 · 측정 책임자 실명 · 범위 밖 전제 · 승인 전 항목 · 반려된 견적 — 이 공급자로서 무엇을 만들어 팔 것인가의 재료가 된다}
\note{교재 "Product 트랙 예고" · 워크북 §7. 혼합 기수이므로 Product 트랙만 듣는 수강생이 아니라 PM 트랙 수강생에게 말하는 프레임이다 — 4일 동안의 발주자 경험이 Day5부터 어떻게 뒤집히는지를 예고한다. 핵심은 정보의 비대칭이다: 발주자였을 때 당신은 이관 목록 6번이 무엇인지 알았지만, 공급자로서는 그것을 밖에서 발견해야 한다. 발주자의 사슬을 끝까지 따라가 본 사람만이 공급자로서 그 사슬의 어느 고리가 비어 있는지를 볼 수 있다. Product 트랙의 교재·도구 명칭은 여기서 말하지 않는다 — Day5 도입에서 한다. 예상 소요 3분.}
\end{frame}

\begin{frame}{이관 목록의 한 줄이 출발점 — "승인 전" 네 줄 중 어느 줄이 누군가의 사업이 되는가}
\begin{tbl}[\scriptsize]{L{0.5cm}L{5.6cm}L{1.6cm}L{1.9cm}L{1.7cm}}
\# & 내용 & 소유자 & 기한 & 무엇이 승인돼야 \\ \midrule
3 & 가맹 발주 마감 22:00 운용 전환 — 물류 야간조 편성과 함께 CAB 재상정 & 한동석 & 2027-09 CAB & 야간조 편성안 · 3PL 시각 합의 \\
6 & 이천센터·안성 검사기록은 P1 범위 밖(종이·엑셀). 98.5\% 목표는 디지털화가 전제 — 2025-11 이사회 반려(34억) 재검토 & 나영선 & 2027-09 이사회 & 이사회 재상정 \\
7 & SAP ECC 2027-12 유지보수 종료 — 애드온·IF-SAP 8본 마이그레이션 계획(KE-04) & 박세연 & 2027-08-31 & 별도 과제 예산 \\
10 & 하자보수기 PMO 잔류 1명 — 유보 3.0 중 1.05 집행 & 프로그램 PMO장 & 2027-06-01 운영위 & 스폰서 결정 \\
\end{tbl}
\keymsg{프로젝트가 끝난 뒤에도 남아 작동하는 문서는 그 프로젝트의 것만이 아니다 — 다음 프로젝트의, 다음 사람의 출발점이다. Day5에 만나자}
\note{워크북 §3.5 항목 6(이관 10건 중 승인 전 4)과 §7의 마지막 질문. 네 줄을 읽고 수강생에게 1분 생각할 시간을 준다 — 어느 줄이 "회사 안의 후속 과제"로 끝나고, 어느 줄이 "밖의 누군가가 만들어 팔 물건"이 되는가. 답을 정해 주지 않는다. 6번(그림 4-4 강조 · 그림 5-3의 분모 B · BR-05)이 편익 98.5\%의 전제이면서 34억 견적이 반려된 자리라는 점만 짚는다 — 반려된 견적은 문제가 사라졌다는 뜻이 아니라 그 값으로는 안 산다는 뜻이다. 예상 소요 3분.}
\end{frame}

\begin{frame}{사후 읽기 · 30일 과제 — 당신 회사의 종료 보고서 한 부로}
\begin{itemize}
\item \textbf{사후 읽기(2주)} — PRINCE2 7 Closing a Project · SRE ch.3\til{}4 · 개정법 조문(다시) · Duffield \& Whitty
\item \textbf{과제 ① 세 칸 O/X} — 원가는 어느 층과 비교됐나 / 이관 목록에 소유자 실명·승인 여부가 있나 / 측정 책임자가 실명인가
\item \textbf{과제 ② 거꾸로 추적} — 그 보고서의 편익 한 줄을 BC·헌장까지 거슬러 고리가 비는 지점 하나를 적는다
\item \textbf{과제 ③ 한 장} — 비는 고리 하나를 채우는 개정 요청(교훈 등록부 ⑤ 소유자 · Syllk 층)을 A4 한 장으로
\item 제출: Day8 종료 전 · 형식은 워크북 §4·§5의 양식 그대로
\end{itemize}
\keymsg{항목의 수가 아니라 세 칸이 문제다 — 30일 뒤 당신 회사의 종료 보고서에 그 세 칸이 생기면 이 4일은 끝난 것이다}
\note{교재 6.9의 과정 전 과제(세 칸 O/X)를 사후 과제로 되돌린다 — 과정 전에는 "읽고 O/X만", 사후에는 "비는 고리를 채우는 개정 요청 한 장"까지. 과제 ②는 그림 6-1의 방법을 자기 회사 문서에 적용하는 것이며, 고리가 비는 지점은 대개 측정 책임자(부서명)·분모(미정의)·승인 전 항목(승인처럼 적힘) 셋 중 하나다. 과제 ③은 워크북 §4 교훈 등록부의 ④ 적용 대상·⑤ 소유자 칸과 §4의 개정 요청 양식을 그대로 쓴다 — 양식이 있어야 process·infrastructure 층(Syllk)에 들어간다. 제출 시점을 Day8 종료 전으로 두어 Product 트랙 수강 중에도 PM 트랙의 문서가 살아 있게 한다. 예상 소요 3분. 이로써 마무리 30분 — 2+4+4+2+5+3+3+3+3+3 = 32분, 필요시 읽을 자료 프레임을 1분으로 줄인다.}
\end{frame}
